% !TEX program = xelatex
% =============================================================================
% Modeling the Probability of Corporate Default
% Using Econometric and Machine Learning Methods
% ICEF HSE formatting: Times New Roman 12pt, 1.5 spacing, A4
%   left 35mm, right 15mm, top/bottom 20mm
% COMPILE WITH XeLaTeX (needed for the Cyrillic title page and annotation).
% =============================================================================

\documentclass[12pt,a4paper]{article}

% -------- fonts and languages (XeLaTeX) --------
\usepackage{fontspec}
\usepackage{polyglossia}
\setmainlanguage{english}
\setotherlanguage{russian}
% Times New Roman covers both Latin and Cyrillic. If it is not installed
% (e.g. on some Overleaf images), replace the three font lines below with a
% metric-compatible open substitute, e.g. "Liberation Serif".
\setmainfont{Times New Roman}
\newfontfamily\cyrillicfont{Times New Roman}
\newfontfamily\cyrillicfonttt{Times New Roman}

\usepackage[a4paper,
            left=35mm,
            right=15mm,
            top=20mm,
            bottom=20mm]{geometry}
\usepackage{amsmath,amssymb,amsthm}
\usepackage{booktabs}
\usepackage{graphicx}
\usepackage{array}
\usepackage{tabularx}
\usepackage{multirow}
\usepackage{caption}
\usepackage{float}
\usepackage[section]{placeins}
\usepackage[round]{natbib}
\usepackage{setspace}
\usepackage{listings}         % for the code appendix
\usepackage{xcolor}
\usepackage{tocloft}          % for table of contents formatting
\usepackage{titlesec}         % section heading control
\usepackage{hyperref}
\hypersetup{
  colorlinks=true,
  linkcolor=black,
  citecolor=black,
  urlcolor=blue,
}

% -------- ICEF spacing and paragraph --------
\onehalfspacing
\setlength{\parskip}{4pt}
\setlength{\parindent}{16pt}

% -------- caption format: "Table N. Title" / "Figure N. Title" --------
\captionsetup[table]{
  labelsep=period,
  name=Table,
  position=above,
  justification=raggedright,
  singlelinecheck=false,
}
\captionsetup[figure]{
  labelsep=period,
  name=Figure,
  position=below,
  justification=centering,
}

% -------- listings style for the Python code appendix --------
\definecolor{codegreen}{rgb}{0.0,0.5,0.0}
\definecolor{codegray}{rgb}{0.5,0.5,0.5}
\definecolor{codepurple}{rgb}{0.55,0.0,0.55}
\definecolor{codeblue}{rgb}{0.0,0.0,0.65}
\lstdefinestyle{python}{
  language=Python,
  basicstyle=\ttfamily\footnotesize,
  keywordstyle=\color{codeblue}\bfseries,
  commentstyle=\color{codegreen}\itshape,
  stringstyle=\color{codepurple},
  numberstyle=\tiny\color{codegray},
  numbers=left,
  numbersep=8pt,
  showstringspaces=false,
  breaklines=true,
  breakatwhitespace=true,
  frame=single,
  rulecolor=\color{codegray},
  columns=fullflexible,
  keepspaces=true,
  tabsize=4,
}

\begin{document}

% ============================================================
%  TITLE PAGE  (ICEF Annex 8 template, Russian)
% ============================================================
\begin{titlepage}
\begin{center}
\onehalfspacing
{\small\textrussian{%
ФЕДЕРАЛЬНОЕ ГОСУДАРСТВЕННОЕ АВТОНОМНОЕ\\
ОБРАЗОВАТЕЛЬНОЕ УЧРЕЖДЕНИЕ ВЫСШЕГО ОБРАЗОВАНИЯ\\
«НАЦИОНАЛЬНЫЙ ИССЛЕДОВАТЕЛЬСКИЙ УНИВЕРСИТЕТ\\
«ВЫСШАЯ ШКОЛА ЭКОНОМИКИ»\\[10pt]
Международный институт экономики и финансов}\\
}
\vfill
{\large\textrussian{Тагиров Ильяс Ильдарович}}\\[26pt]
{\Large\bfseries\textrussian{МОДЕЛИРОВАНИЕ ВЕРОЯТНОСТИ\\[4pt]
КОРПОРАТИВНОГО ДЕФОЛТА\\[4pt]
С ИСПОЛЬЗОВАНИЕМ МЕТОДОВ ЭКОНОМЕТРИКИ И\\[4pt]
МАШИННОГО ОБУЧЕНИЯ}}\\[14pt]
{\large (Modeling the Probability of Corporate Default Using\\
Econometric and Machine Learning Methods)}\\[34pt]
{\normalsize\textrussian{%
Курсовая работа\\[6pt]
по направлению подготовки 38.03.01 «Экономика»\\
образовательная программа «Международная программа\\
по экономике и финансам»}\\
}
\vfill
{\normalsize\textrussian{Москва, 2026}}\\[18pt]
\end{center}
\begin{flushright}
{\normalsize\textrussian{%
Научный руководитель\\
Каныгина Надежда Геннадьевна}\\[20pt]
\rule{60mm}{0.4pt}\\
\hspace{6mm}подпись \hspace{26mm} дата
}
\end{flushright}
\end{titlepage}

% ============================================================
%  ABSTRACT (English) + ANNOTATION (Russian) + TABLE OF CONTENTS
%  ICEF order: title page -> abstract -> annotation -> contents
% ============================================================
\newpage
\begin{abstract}
\noindent
This paper compares classical econometric bankruptcy models with modern
machine-learning classifiers on the Polish Companies Bankruptcy dataset of
\cite{zieba2016}. The main empirical finding is that missing financial-ratio
values are not merely a nuisance to be imputed away: the pattern of
non-disclosure is itself a strong distress signal, appearing among the largest
LASSO coefficients and among the most important XGBoost features. We therefore treat missingness as an informative component of corporate
default prediction, grounding the approach in Rubin's missing-data taxonomy
and the missing-indicator method rather than claiming it as new
\citep{rubin1976,jones1996}. The second contribution is to separate the
machine-learning advantage into two parts: access to a broader 75-variable
information set, including missingness flags, and the additional nonlinear
modeling power of tree ensembles. The third contribution is a unified,
leakage-free benchmark. We re-estimate Altman, Ohlson, and Zmijewski-style
models and compare them with LASSO, random forest, and XGBoost using the same
stratified folds, fold-internal imputation, class-weighted training, and the
same discrimination, calibration, and cost metrics. The results show that
XGBoost dominates the classical models, but also that a large part of this gap
comes from feature breadth and informative missingness, not from algorithmic
flexibility alone.
\end{abstract}

\begin{otherlanguage}{russian}
\begin{center}
\textbf{Аннотация}
\end{center}
\noindent
В работе сравниваются классические эконометрические модели банкротства и
современные классификаторы машинного обучения на наборе данных о банкротстве
польских компаний \cite{zieba2016}. Основной эмпирический результат состоит в
том, что пропущенные значения финансовых коэффициентов не являются просто
помехой, которую нужно восполнить: сам факт отсутствия отчётности является
сильным сигналом неблагополучия и проявляется как среди наибольших
коэффициентов LASSO, так и среди наиболее важных признаков XGBoost. Поэтому мы
рассматриваем пропуски как информативную составляющую прогнозирования дефолта,
опираясь на классификацию пропущенных данных Рубина и метод индикаторов
пропусков, не претендуя на новизну самого метода \citep{rubin1976,jones1996}.
Второй вклад работы состоит в разложении преимущества машинного обучения на две
части: доступ к более широкому набору из 75 переменных, включая индикаторы
пропусков, и дополнительную нелинейную выразительность древесных ансамблей.
Третий вклад — единый и защищённый от утечки протокол сравнения. Мы заново
оцениваем модели типа Альтмана, Ольсона и Змиевски и сравниваем их с LASSO,
случайным лесом и XGBoost на одних и тех же стратифицированных разбиениях, с
восполнением пропусков внутри каждого разбиения, с взвешиванием классов и с
одними и теми же метриками различения, калибровки и стоимости ошибок.
Результаты показывают, что XGBoost превосходит классические модели, но большая
часть этого разрыва объясняется широтой набора признаков и информативностью
пропусков, а не только гибкостью алгоритма.
\end{otherlanguage}

\clearpage
\tableofcontents
\clearpage
% =============================================================================
\section{Introduction}
% =============================================================================

When a company goes bankrupt, the consequences spread widely. Creditors lose
money, suppliers lose customers, employees lose jobs, and investors absorb
losses. When many firms fail at the same time, as happened during the
2007-2009 financial crisis, the effects ripple through the broader economy.
This is why predicting corporate bankruptcy has been an active research area
for over fifty years, attracting work from academics, regulators, and
practitioners alike
\citep{altman1968, ohlson1980, zmijewski1984, shumway2001, campbell2008,
barboza2017, zieba2016, mai2019}.

\subsection{Two strands of research}

The literature falls into two broad approaches. The first is
\textbf{econometric}: it uses financial ratios from company accounts and
estimates formal statistical models with interpretable coefficients. The key
early contributions are \cite{altman1968}, who used five accounting ratios
and a technique called multiple discriminant analysis to classify 66 U.S.\
manufacturing firms as bankrupt or healthy; \cite{ohlson1980}, who moved to
logistic regression with nine variables and a much larger sample of 2,163
firms (105 bankrupt, 2,058 healthy) over 1970-1976; and
\cite{zmijewski1984}, who used a probit model with three ratios on 840 firms
over 1972-1978. \cite{shumway2001} extended this tradition by treating
bankruptcy as a time process using a hazard model, which allows the
predictors to change over time and correctly handles firms that are observed
for different numbers of years. These models are valued for their
transparency: a coefficient tells you directly how a change in, say,
retained earnings relative to total assets affects the probability of failure.

The second approach is \textbf{machine learning} (ML). Starting in the
late 1980s with neural networks \citep{odom1990} and expanding to support
vector machines \citep{shin2005, min2005}, random forests \citep{barboza2017},
gradient-boosted trees \citep{zieba2016, carmona2019}, and more recently
deep learning \citep{mai2019, pellegrino2024}, ML methods have been applied
to bankruptcy prediction with growing frequency. The common finding across
these studies is that ML models achieve higher out-of-sample discrimination
than classical econometric benchmarks, as measured by the area under the
ROC curve (AUC).

\subsection{A problem with existing comparisons}

The claim that ``ML beats logit'' is now widespread. However, looking closely
at how existing studies conduct their comparisons reveals that the ML and
econometric models are usually not tested under the same conditions. Three
differences appear regularly.

\textbf{Imbalance handling.} Bankruptcy is a rare event. In most datasets,
fewer than 5\% of firms go bankrupt in any given year, and often fewer than
1\% \citep{shumway2001}. Many ML studies deal with this imbalance by
resampling the training data, either by randomly removing healthy firms
until the classes are roughly equal (undersampling, as in
\citealt{pellegrino2022}), or by generating synthetic bankrupt firms via
SMOTE \citep{chawla2002}. The econometric baseline, meanwhile, is usually
fitted on the original imbalanced data. This creates a hidden problem: a model trained on 50/50 balanced data will
produce predicted probabilities that are too high when applied to a real-world
dataset where the default rate is around 4\% \citep{king2001, saerens2002}.
The model may rank firms well but will be poorly calibrated.

\textbf{Feature set.} ML models are often given many more input variables
than the econometric baseline. An ML model with 64 features that outperforms
a 5-variable Altman logit is providing evidence about the value of those
extra features, not necessarily about ML algorithms being intrinsically
superior.

\textbf{Evaluation metric.} Most studies report only the AUC, which measures
how well a model \emph{ranks} firms by risk. A model can rank firms well but
still produce badly wrong probability estimates. Calibration metrics, which
test whether a model's predicted probability of 10\% actually corresponds to
a 10\% observed default rate, capture a different and equally important
dimension of performance. Reporting only AUC hides this.

\subsection{This paper's contribution}

We revisit the ML-vs-econometric question using the Polish Companies
Bankruptcy dataset of \cite{zieba2016}. The contribution is not simply that we
run several models and compare their AUCs. That would be a useful benchmark,
but it would not be enough as a research contribution. The paper's contribution
is more specific and rests on three claims.

\begin{enumerate}

\item \textbf{Informative missingness as a distress signal.} We show that the
pattern of missing financial-ratio values predicts bankruptcy and should be
interpreted economically, not treated only as a preprocessing inconvenience.
This is the strongest empirical result in the paper: firms close to bankruptcy
are much more likely to have missing ratio values, the missingness indicators
receive some of the largest LASSO coefficients, and they also appear among the
important XGBoost features. Methodologically, this is an application of the
missing-indicator method and of the idea that data may be missing not at random
(MNAR) \citep{rubin1976,jones1996,little2019}. The novelty claimed here is
therefore not the indicator technique itself, but its use and interpretation as
a corporate-distress signal in this bankruptcy dataset.

\item \textbf{Disentangling feature breadth from algorithmic power.} Existing
ML-vs-econometric comparisons often give the ML model a much richer feature set
than the classical benchmark, and then attribute the entire gain to the
algorithm. We separate these two effects. A logit estimated on a broad feature
set tests how much is gained from additional accounting information and
missingness flags while keeping the model linear. XGBoost on the same broad
feature set then tests how much is gained from nonlinear thresholds and
interactions. This matters because \cite{barboza2017} also show that adding
complementary financial indicators materially improves default prediction; our
analysis makes that feature-breadth channel explicit on the Polish dataset.

\item \textbf{A unified and leakage-free evaluation protocol.} We apply one
protocol to all models: stratified five-fold cross-validation, training-fold
median imputation, class-weighted loss functions rather than pre-split
resampling, and the same set of discrimination, calibration, and economic-cost
metrics. This directly addresses a common weakness in the literature, where ML
models and econometric baselines are often evaluated under different feature
sets, different imbalance treatments, or different metrics \citep{hand2009,
king2001}. The aim is not to give ML a favorable comparison,
but to make the comparison hard to dismiss.

\end{enumerate}

The paper proceeds as follows. Section~\ref{sec:lit} reviews the
accounting-based and ML literatures and explains the methodological problems
with existing comparisons. Section~\ref{sec:data} describes the data.
Section~\ref{sec:methodology} specifies the models and evaluation protocol.
Section~\ref{sec:results-econ} reports the re-estimated econometric baselines
and Section~\ref{sec:results-ml} the machine-learning models and the full
comparison. Section~\ref{sec:figures} presents the figures and the
Brier-score decomposition, Section~\ref{sec:horizons} examines performance
across prediction horizons, and Section~\ref{sec:shap} interprets the XGBoost
model with SHAP values. Section~\ref{sec:robustness} reports robustness
checks and Section~\ref{sec:conclusion} concludes.

\clearpage
% =============================================================================
\section{Literature Review}\label{sec:lit}
% =============================================================================

We organize the literature into five parts: classical accounting-based models,
hazard and panel models, machine-learning models, evaluation methodology, and
missing data in distress prediction. The paper sits at the intersection of all
five: it compares prediction models, asks whether the comparison is fair,
considers what a proper panel design would look like, and treats missingness
as a signal rather than a nuisance.

\subsection{Classical accounting-based models}

\paragraph{Altman (1968): multiple discriminant analysis.}
\cite{altman1968} is the founding paper of quantitative bankruptcy prediction.
Altman worked with 66 publicly traded U.S.\ manufacturing firms, 33 that had
filed for bankruptcy under Chapter~X of the National Bankruptcy Act between
1946 and 1965, matched to 33 healthy firms of similar size and industry.
He applied \textit{multiple discriminant analysis} (MDA), a technique from
\cite{fisher1936} that finds the linear combination of variables that best
separates two groups, assuming they follow a joint normal distribution with
equal covariances. Starting from 22 candidate ratios, he selected five. The resulting \textit{Z-score} is:
\begin{equation}
Z = 1.2 \cdot X_1 + 1.4 \cdot X_2 + 3.3 \cdot X_3 + 0.6 \cdot X_4 + 1.0 \cdot X_5
\label{eq:altman}
\end{equation}
where $X_1$ is working capital over total assets, $X_2$ is retained earnings
over total assets, $X_3$ is EBIT over total assets, $X_4$ is the market value
of equity over total liabilities, and $X_5$ is sales over total assets. A firm
with $Z < 1.81$ was classified as likely to fail; above $2.99$ as likely to
survive; the range in between was called the ``zone of ignorance.'' Altman
reported 95\% accuracy one year before bankruptcy and 72\% two years before,
on his matched sample.

The Z-score has been updated several times, including versions for private and
non-manufacturing firms \citep{altman2000}, and remains widely used in
professional credit analysis. Its main weaknesses are the normality and equal-covariance assumptions,
the distortion from matched sampling when the real bankruptcy rate is far
below 50\%, and the use of market capitalization, which is unavailable
for private firms.

\paragraph{Ohlson (1980): logistic regression.}
\cite{ohlson1980} addressed the limitations of MDA by switching to logistic
regression. Unlike MDA, the logit model makes no distributional assumption
and directly estimates the probability that a firm goes bankrupt:
\begin{equation}
\Pr(y_i = 1 \mid \mathbf{x}_i) =
\frac{1}{1 + \exp(-\mathbf{x}_i^\top \boldsymbol{\beta})},
\label{eq:logit}
\end{equation}
where $y_i = 1$ indicates bankruptcy, $\mathbf{x}_i$ is the vector of
predictors, and $\boldsymbol{\beta}$ is estimated by maximum likelihood.
Ohlson used 105 bankrupt and 2,058 non-bankrupt U.S.\ industrial firms from
Compustat over 1970-1976. His nine predictors included size, leverage,
working capital, liquidity, profitability, cash flow, a negative-book-equity
indicator, and two variables requiring prior-year information. These improvements over Altman, a larger sample, no distributional assumption,
and direct probability output, made the O-score the standard academic benchmark.

\paragraph{Zmijewski (1984): probit regression.}
\cite{zmijewski1984} proposed a three-variable probit model. The probit is
closely related to the logit, but uses the standard normal cumulative
distribution function rather than the logistic CDF:
\begin{equation}
\Pr(y_i = 1 \mid \mathbf{x}_i) =
\Phi\!\left(\mathbf{x}_i^\top \boldsymbol{\beta}\right),
\label{eq:probit}
\end{equation}
where $\Phi(\cdot)$ is the standard normal CDF. Zmijewski's predictors are net
income over total assets, total liabilities over total assets, and current
assets over current liabilities. The model is simple but important because it explicitly addresses choice-based
sampling bias: when bankrupt firms are overrepresented in the sample, the
intercept is biased but the slope coefficients remain consistent. Later work
finds that no single specification consistently outperforms the others, which
is why we evaluate all three side by side \citep{mossman1998}.

\subsection{Hazard and panel models}

\cite{shumway2001} argued that single-period models are statistically
inconsistent for bankruptcy prediction because firms are observed over time
and may survive, fail, or exit the sample. The right framework is a
\textit{discrete-time hazard model}, which estimates the probability of
failure in period $t$ conditional on having survived to $t$. Shumway showed
that this can be run as a pooled logit on firm-year observations, with
time-varying covariates. On U.S.\ data, several classical accounting ratios
lost significance once the hazard structure was imposed, while market variables
such as firm size, past returns, and idiosyncratic volatility became strong
predictors.

\cite{chava2004} extend this logic by adding industry effects and show that
industry differences affect both baseline failure risk and the slopes on key
variables. \cite{campbell2008} develop a dynamic logit model using both
accounting and market information, including leverage, profitability, market
capitalization, past returns, volatility, cash, market-to-book, and price per
share, and show that these variables capture much of the time variation in
aggregate failure rates. The implication for our paper is direct: a hazard model is the modern
econometric benchmark. It is not feasible here because the Polish dataset is
anonymized and provides no firm identifiers or calendar-year stamps, as
explained in Section~\ref{sec:panel-limits}.

\subsection{Machine-learning models for bankruptcy prediction}

ML methods were first applied to bankruptcy prediction in the late 1980s with
simple neural networks \citep{odom1990}. The field has since expanded to
include further neural-network distress models \citep{altman1994}, support
vector machines \citep{shin2005, min2005}, tree ensembles
\citep{barboza2017,jones2017}, gradient boosting \citep{zieba2016,
carmona2019}, and recurrent neural networks \citep{mai2019,pellegrino2024}.
Recent surveys confirm that this literature has grown quickly. \cite{dasilas2024}
review 207 empirical bankruptcy-prediction studies from 2012 to mid-2023,
covering the rapid shift toward ensemble and hybrid machine-learning methods.

\cite{barboza2017} compared several ML methods, including support vector
machines, random forests, bagging, boosting, and neural networks, against
logistic regression and discriminant analysis on a large sample of North
American firms. They found that ML methods were roughly 10 percentage points
more accurate than traditional models. Crucially for this paper, they also
show that adding complementary financial indicators to the classic Altman
variables materially raises performance. That finding supports our argument
that part of the ML advantage is really a feature-breadth advantage.

\cite{jones2017} evaluate 16 statistical and machine-learning classifiers and
find that newer ensemble learners such as generalized boosting, AdaBoost, and
random forests perform best, while simple logit and discriminant models still
remain useful benchmarks. \cite{zieba2016}, whose dataset we use, applied
XGBoost to Polish firms observed between 2000 and 2013 and reported AUC values
of 0.93-0.96 across the five prediction horizons. Their result is the starting
point for our paper, but we differ in focus: we compare econometric and ML
models under one protocol and then separate feature breadth from nonlinear
modeling power.

\cite{pellegrino2022} applied deep neural networks and tree ensembles to U.S.\
public companies over 1999-2018, with the best model reaching AUC 0.85.
However, their training protocol balanced the sample to 50/50 by dropping
healthy firms. As \cite{king2001} showed, this inflates predicted probabilities
unless the intercept is corrected. We treat such results as evidence that ML
can help, but also that protocol consistency matters.

\subsection{Evaluation methodology: imbalance, leakage, and metrics}

Bankruptcy prediction is inherently an imbalanced classification problem, and
a body of statistical work warns against ignoring this. \cite{saito2015}
showed that PR-AUC is more informative than ROC-AUC when the minority class is
rare. ROC-AUC measures ranking over the whole sample, while PR-AUC directly
asks whether the firms flagged as risky are actually bankrupt. \cite{hand2009}
argues that even AUC itself is not always a coherent measure of classifier
performance, which supports reporting a battery of metrics rather than one
headline number.

\cite{king2001} documented how balanced-sample training inflates absolute
probability estimates. \cite{niculescu2005} showed that even strong ranking
models such as random forests can produce poorly calibrated probabilities.
These issues are why we report both discrimination and calibration metrics,
including the Brier score, reliability diagrams, expected calibration error,
and post-hoc isotonic recalibration.

A separate problem is data leakage. Resampling methods such as SMOTE
\citep{chawla2002} must be applied inside the training fold only. Applying
SMOTE before the train/test split lets synthetic observations carry information
from the validation data into training, inflating reported performance. The
same logic applies to imputation and scaling, which must be fit on training
data alone. Our evaluation protocol therefore fits
imputation, scaling, and any resampling operation inside each cross-validation
fold only.

\subsection{Missing data in distress prediction}

The strongest novel element of this paper is not that we impute missing values,
but that we treat the missingness pattern as economically informative. The
statistical foundation is Rubin's taxonomy of missing completely at random,
missing at random, and missing not at random \citep{rubin1976}, developed
further in \cite{little2019}. In corporate bankruptcy data, missingness can
plausibly be MNAR: the value may be absent precisely because the firm is in
financial trouble, because the denominator of a ratio is zero or very small,
or because reporting discipline deteriorates as the firm approaches failure.

The missing-indicator method, adding a binary flag for each missing value,
is a long-standing technique \citep{jones1996}. In prediction it is useful
because missing data can carry information that would be discarded if the
row were simply dropped. \cite{granularimputation2024} apply this idea
directly to the Polish dataset and confirm that missingness patterns matter.
We use missingness indicators as a predictive feature, while being clear
that the method itself is not new.

\subsection{Where this paper fits}

The gap in the literature is specific. A large body of empirical work has
established that ML methods produce higher out-of-sample AUC than classical
accounting-ratio models \citep{barboza2017, jones2017, dasilas2024}. That
finding is not in dispute here. What is disputed is whether the reported gaps
are as large and robust as they appear, and what is actually driving them.

Three problems recur in the comparison literature.
First, ML models are typically given access to many more features than the
econometric baseline; the performance gap therefore reflects both algorithmic
power and information advantage, and most studies do not separate these.
Second, imbalance is often handled differently across models: some studies
balance the training sample for ML models but not for the logit baseline,
or apply SMOTE before the train/test split, inflating the reported ML
performance \citep{king2001, chawla2002}. Third, most studies report only
ROC-AUC, which is relatively insensitive to the handling of the rare positive
class; PR-AUC and calibration metrics tell a more complete story
\citep{saito2015, hand2009}.

The Polish dataset is a natural setting for addressing these problems because
it is the most widely used public benchmark for this literature, covering more
than 40,000 firm-horizon observations from an emerging market with significant
variation in macro conditions. \cite{zieba2016} used it to introduce XGBoost
to bankruptcy prediction and reported AUC values of 0.93-0.96, but their
study does not include re-estimated classical benchmarks, does not apply a
unified evaluation protocol, and does not separate feature breadth from
algorithmic flexibility. No subsequent study has filled all three gaps
simultaneously on this dataset. This paper does, and in doing so offers a
more precise account of where the ML advantage comes from and how much of it
survives a controlled comparison.

\clearpage
% =============================================================================
\section{Data}\label{sec:data}
% =============================================================================

\subsection{Source}

We use the \textit{Polish Companies Bankruptcy} dataset originally compiled by
\cite{zieba2016} and publicly available through the UCI Machine Learning
Repository (dataset ID~365).\footnote{\url{https://doi.org/10.24432/C5F600};
CC~BY~4.0 licence.} The underlying data come from the Emerging Markets
Information Service (EMIS), a commercial database of corporate filings in
emerging markets. Healthy firms are observed over fiscal years 2007-2013;
bankrupt firms are observed in the years preceding their bankruptcy filings,
which occurred between 2000 and 2012.

The dataset provides 64 pre-computed financial ratios for each firm-year
observation, covering liquidity, profitability, leverage, activity, and
efficiency. Importantly, these ratios were computed by the original authors
directly from the underlying financial statements, rather than being left for
the researcher to construct. This avoids a common source of error in datasets
that expose only raw accounting items, where the units of reported figures can
differ across firms and years and produce meaningless ratios when divided
directly.

\subsection{Horizon structure}
\label{sec:horizon-structure}

The dataset is distributed as five separate files, one for each prediction
horizon. The original file names are counterintuitive: \texttt{1year.arff} contains
statements from the \emph{earliest} year of the forecasting window, while
\texttt{5year.arff} contains statements from the \emph{latest} year,
predicting bankruptcy within one year. To avoid confusion, we rename
the files using the prediction horizon length:

\begin{center}
\begin{tabular}{lll}
\toprule
UCI file & Our label & Prediction horizon \\
\midrule
\texttt{5year.arff} & H1 & 1 year ahead  \\
\texttt{4year.arff} & H2 & 2 years ahead  \\
\texttt{3year.arff} & H3 & 3 years ahead  \\
\texttt{2year.arff} & H4 & 4 years ahead  \\
\texttt{1year.arff} & H5 & 5 years ahead  \\
\bottomrule
\end{tabular}
\end{center}

\noindent H1 is therefore the shortest and most informative prediction task
(financial statements are close to the bankruptcy event), and H5 is the
longest and least informative. H1 is our primary analysis horizon; H2-H5
appear in the horizon-degradation analysis of Section~5.

\subsection{Macroeconomic context of the observation windows}
\label{sec:macro-context}

The dataset provides no calendar year for individual observations, so
macroeconomic variables cannot be merged at the firm-year level. Attaching,
for example, GDP growth for a specific year to a specific row would be
false precision. What we can do is describe the macro environment covered
by the data and note how it may affect the baseline bankruptcy rate.

Table~\ref{tab:polish-macro} reports the Polish macro series for 2000-2013,
the period covered by the bankrupt and still-operating firms in the dataset.
The window includes Poland's pre-crisis expansion, the 2008-2009 global crisis,
and the slower post-crisis period. Poland is unusual in this period because it
continued to record positive real GDP growth in 2009, while the EU as a whole
fell into recession. This resilience does not remove macro risk from the data,
but it means that the Polish default environment is not directly comparable to
countries where the crisis produced a deep output collapse. Because individual
observations carry no calendar date, macro variables cannot be merged at the
firm-year level. They can, however, be attached at the level of the horizon
window, since each of the five files corresponds to a different distance from
the outcome and therefore to a different average position in the macro cycle.
We use this horizon-level macro information in the pooled analysis of
Section~\ref{sec:pooled}, where we also show that it is largely absorbed by
the horizon dummies.

\begin{table}[H]
\centering
\caption{Polish macroeconomic context, 2000--2013.\textsuperscript{*}}
\label{tab:polish-macro}
\small
\resizebox{0.92\textwidth}{!}{%
\begin{tabular}{lrrrrrrrrrrrrrr}
\toprule
 & 2000 & 2001 & 2002 & 2003 & 2004 & 2005 & 2006 & 2007 & 2008 & 2009 & 2010 & 2011 & 2012 & 2013 \\
\midrule
Real GDP growth & 4.66 & 1.23 & 1.90 & 3.52 & 5.09 & 3.26 & 6.20 & 6.76 & 4.38 & 2.62 & 3.17 & 5.26 & 1.51 & 0.68 \\
CPI inflation & 10.1 & 5.5 & 1.9 & 0.7 & 3.5 & 2.2 & 1.2 & 2.5 & 4.3 & 3.8 & 2.6 & 4.2 & 3.7 & 1.1 \\
Unemployment & 15.1 & 17.5 & 20.0 & 20.0 & 19.0 & 17.6 & 14.8 & 11.2 & 9.5 & 12.1 & 12.4 & 12.5 & 13.4 & 13.4 \\
NBP reference rate & 19.0 & 11.5 & 6.75 & 5.25 & 6.5 & 4.5 & 4.0 & 5.0 & 5.0 & 3.5 & 3.5 & 4.5 & 4.25 & 2.5 \\
\bottomrule
\end{tabular}%
}
\caption*{\footnotesize \textsuperscript{*} GDP growth and CPI inflation are annual
percentages; unemployment is the registered year-end rate; the NBP reference
rate is the year-end policy rate. Sources: World Bank national accounts, GUS
(registered unemployment), National Bank of Poland.}
\end{table}

\subsection{Sample sizes and class imbalance}

Table~\ref{tab:horizons} shows the sample sizes, bankruptcy counts, and
default rates by horizon. Three features stand out.

First, the bankruptcy rate is low across all horizons, ranging from 6.9\%
at H1 to 3.9\% at H5. This reflects the fact that the shorter-horizon files
contain a higher proportion of firms that are already in clear distress,
while the longer-horizon files include many firms that will eventually fail
but have not yet entered that phase.

Second, all five horizons have a severely imbalanced class ratio, roughly
14 healthy firms for every bankrupt one. This is the core imbalance problem
that shapes our entire evaluation strategy.

Third, the raw number of bankrupt firms is small: only 271 at H5 and 410 at
H1. This matters because the effective statistical signal available to the
models is driven by the minority class. A model trained on 271 bankrupt
observations cannot reliably learn subtle nonlinear patterns; this should
temper expectations about the practical benefit of adding model complexity.

\begin{table}[ht]
\centering
\caption{Sample sizes and class balance by prediction horizon.\textsuperscript{*}}
\label{tab:horizons}
\centering
\resizebox{\textwidth}{!}{%
\begin{tabular}{lrrrrrr}
\toprule
 & Obs. & Bankrupt & Rate (\%) & Complete rows & BR complete (\%) & BR incomplete (\%) \\
\midrule
H1 (1 yr ahead) & 5{,}910  & 410 & 6.94 & 3{,}031 & 3.37 & 10.70 \\
H2 (2 yr ahead) & 9{,}792  & 515 & 5.26 & 4{,}769 & 2.52 &  7.86 \\
H3 (3 yr ahead) & 10{,}503 & 495 & 4.71 & 4{,}885 & 2.19 &  6.91 \\
H4 (4 yr ahead) & 10{,}173 & 400 & 3.93 & 4{,}088 & 1.79 &  5.37 \\
H5 (5 yr ahead) & 7{,}027  & 271 & 3.86 & 3{,}194 & 0.94 &  6.29 \\
\bottomrule
\end{tabular}%
}
\caption*{\footnotesize \textsuperscript{*} The final two columns show the bankruptcy rate
(BR) separately for complete and incomplete rows, documenting the informative
missingness discussed in Section~\ref{sec:missing}. Source: author's
calculations.}
\end{table}

\subsection{The 64 financial ratios}
\label{sec:ratios}

Each observation in the dataset contains 64 pre-computed ratios derived from
the firm's annual financial statements. We rename the original variable codes
(Attr1 to Attr64) to readable labels that reflect their economic meaning
(e.g.\ \texttt{EBIT\_TA} for EBIT divided by total assets, \texttt{TL\_TA}
for total liabilities over total assets).

The 64 ratios include all five components of the \cite{altman1968} Z-score:

\begin{itemize}
\item \textbf{Working capital / Total assets}, Altman's $X_1$,
      our \texttt{WC\_TA} (Attr3).
\item \textbf{Retained earnings / Total assets}, Altman's $X_2$,
      our \texttt{RE\_TA} (Attr6).
\item \textbf{EBIT / Total assets}, Altman's $X_3$,
      our \texttt{EBIT\_TA} (Attr7).
\item \textbf{Book value of equity / Total liabilities}, Altman's $X_4$
      (using book equity in place of market value for compatibility with
      private firms), our \texttt{BE\_TL} (Attr8).
\item \textbf{Sales / Total assets}, Altman's $X_5$,
      our \texttt{sales\_TA} (Attr9).
\end{itemize}

Seven of the nine components of \cite{ohlson1980}'s O-score are also
directly available (SIZE, TLTA, WCTA, CLCA, NITA, FUTL, and OENEG). The
remaining two, INTWO (net income negative in both the current and prior
year) and CHIN (normalized change in net income), require two consecutive
observations for the same firm.\footnote{The dataset is anonymized: firms
are not identified across rows, so it is impossible to link this year's
observation with last year's for the same firm. This also rules out the
\cite{shumway2001} discrete-time hazard model, which requires a panel
structure.} Our Ohlson specification therefore uses seven of the original
nine variables, which we document explicitly in the methodology section.
All three components of \cite{zmijewski1984}'s X-score are present
(NITA, TLTA, and the current ratio).

\subsection{Missing values and what they signal}
\label{sec:missing}

Unlike some datasets where rows with missing values are dropped before
release, the Polish data keep them. For most of the 64 ratios the missingness
rate is negligible. One feature stands out: \texttt{CA\_inv\_LTL} (current
assets minus inventories divided by long-term liabilities) is missing in 43\%
of H1 rows, because firms with no long-term debt produce a division by zero.

The important point is that this missingness is not random with respect to
bankruptcy. Table~\ref{tab:horizons} shows, for each horizon, the bankruptcy
rate among rows with no missing values (``complete rows'') and among rows
with at least one missing value (``incomplete rows''). The gap is
consistently large:

\begin{itemize}
\item At H1, the bankruptcy rate among incomplete rows (10.7\%) is
      3.2 times the rate among complete rows (3.4\%).
\item At H5, the gap is even larger: 6.3\% versus 0.9\%, a factor of 6.7.
\end{itemize}

This makes sense. Distressed firms may have no long-term debt, may
restructure their balance sheet in ways that leave items undefined, or may
simply become less disciplined in their reporting as they approach failure.
A missing value is itself a signal of distress.

Simply replacing missing values with the sample median, without recording
that the value was missing, would throw away this information. We therefore adopt a two-step approach. First, for every
feature that is missing in at least 1\% of rows in any horizon, we add a
binary indicator column (e.g.\ \texttt{miss\_CA\_inv\_LTL}) that equals 1
whenever the original value is missing. This produces 11 such indicator
columns in total. Since these indicators are derived directly from the
pattern of missing values, no model fitting is involved, they do not
introduce any risk of data leakage and can safely be added before the
cross-validation loop. Second, the missing values in the original ratio
columns are filled with the \textit{training-fold median} during
cross-validation, fitted on the training portion of each fold only
(Section~\ref{sec:methodology}).

As a robustness check (Section~7), we repeat the main analysis on the
complete-case subsample (only rows with no missing values). This lets the
reader judge whether our results are driven by the imputation choices or
hold up on the fully observed data.

\subsection{Panel structure: what we can and cannot do}
\label{sec:panel-limits}

A natural econometric extension is a discrete-time hazard model in the spirit
of \cite{shumway2001}. In such a design, the unit of observation is a
firm-year spell. For each firm and each year in which the firm is still alive,
one estimates the probability that the firm fails in that period conditional
on having survived up to the beginning of the period. Operationally, this can
be estimated as a pooled logit with a time variable, period effects, and
possibly time-varying accounting, market, industry, and macro covariates. This
framework handles censoring more naturally than a static cross-section and is
closer to the way bankruptcy risk evolves in reality. More recent multiperiod
approaches, such as the forward-intensity model of \cite{duan2012}, extend the
same idea to predict default probabilities over several future horizons at
once.

The dataset allows part of this idea to be implemented and rules out another
part, and it is important to be clear about which is which.

What we \emph{can} do is exploit the horizon structure of the five files. Each
file corresponds to a different distance between the financial statement and
the bankruptcy outcome, from one year (H1) to five years (H5). Pooling the
files and adding a dummy for each horizon plays the role of the period effects
in a hazard model: the dummies absorb the fact that the baseline default rate
falls as the horizon lengthens. Going further, interacting the financial ratios
with the horizon dummies, the slope-dummy idea, tests whether the same ratio
carries the same meaning at every distance from the event, or whether its
predictive content changes. We implement both of these in
Section~\ref{sec:pooled}, and they produce one of the more interesting findings
in the paper: the predictive power of the standard ratios is concentrated close
to the event and fades as the horizon lengthens.

What we \emph{cannot} do is estimate firm-level effects. The observations are
anonymized, so there is no firm identifier, and there is no published key
linking a firm's H1 observation to its H2, H3, H4, or H5 observation. Without a
firm ID we cannot follow a single firm across years, cannot build a true risk
set, cannot identify individual survival spells, and cannot estimate firm fixed
effects that would absorb time-invariant unobserved heterogeneity such as
business model or reporting quality. There is also no calendar-year stamp at
the row level, so macroeconomic variables can be attached only at the level of
the horizon window, not the individual firm-year. For the same reason, the two
Ohlson variables that require prior-year information, INTWO and CHIN, cannot be
reconstructed.

The practical consequence is a clear division of labour. Horizon dummies, slope
interactions, and horizon-level macro variables are feasible and are used.
Firm fixed effects, firm-level macro merges, and a fully specified hazard model
with market and industry controls would require a linkable panel such as
Compustat, Amadeus, SPARK, or a V4FinBench-style dataset with firm identifiers
and calendar time \citep{kostrzewa2025}. Extending the analysis to such data is
the natural next step and would test whether the missingness effect documented
here survives in a true panel.

Within each horizon file, we use \textit{stratified five-fold
cross-validation} as the primary evaluation framework: observations are
divided into five folds with stratification on the bankruptcy label, so that
each fold contains approximately the same proportion of bankrupt firms.
Stratification matters because with only 271-515 bankrupt observations per
horizon, a purely random split could produce folds with quite different class
rates, adding unnecessary noise to the performance estimates.


\subsection{Data limitations}
\label{sec:data-limits}

The dataset has several properties that set the boundaries of the analysis.
It is useful to state them together.

\paragraph{Anonymization.} Firms are identified only by a row number within
each horizon file. There are no names, registration numbers, or any other
identifier. This means all analysis is confined to the 64 pre-computed ratios provided.

\paragraph{No calendar-year stamps.} Individual observations carry no date.
The documentation states that bankrupt firms span 2000-2012 and
still-operating firms span 2007-2013, but these are aggregate windows for
the whole sample, not labels on individual rows. We cannot assign a
specific year to any single observation, which rules out firm-level macro
merges and firm-year panels.

\paragraph{No market-based variables.} Modern default models such as
\cite{shumway2001} and \cite{campbell2008} show that equity market variables,
including market capitalization, past stock returns, idiosyncratic volatility,
and distance-to-default in the spirit of \cite{merton1974}, carry substantial
predictive information beyond accounting ratios. All firms in this dataset are
private or anonymized, so equity prices are unavailable. This is a genuine limitation: part of what makes the three-variable
econometric models look weak relative to the 75-variable ML model may simply
be that the ML model sees more accounting information.

\paragraph{No industry or firm-age information.} \cite{chava2004} show that
default rates and the marginal impact of financial ratios differ across
industries. Firm age is also a standard control in modern distress models.
Neither is available here. The comparisons therefore use a limited set of controls and should be read
as within-dataset benchmarks, not claims about the full information set.

\paragraph{Potential overlap across horizon files.} The five files are defined
by the length of the forecasting window, not by calendar time. It is possible
that the same firm appears in more than one file, observed at different points
before its eventual bankruptcy or survival. Without firm identifiers we cannot
detect or adjust for this overlap, which means the five files should be treated
as independent cross-sections rather than as panels.

These limitations do not invalidate the analysis. The results show how much
six model classes can extract from the same set of accounting ratios and
missingness indicators, under a protocol that keeps the comparison fair. Extending
the analysis to a dataset with firm identifiers, market data, and industry
codes would be a natural next step; the V4FinBench benchmark of
\cite{kostrzewa2025} is a candidate for this.

\subsection{Descriptive statistics of the Altman ratios}
\label{sec:descriptives}

Before fitting any models, we check whether the five Altman ratios actually
differ between bankrupt and healthy firms in the Polish data, and whether the
direction matches theoretical expectations. Table~\ref{tab:descriptives}
reports group means and standard deviations at H1, together with two tests
of group separation: Welch's two-sample $t$-test and Cohen's $d$.

\textbf{Why Welch's $t$-test?} The standard $t$-test assumes equal
variances across groups. This is violated here: distressed firms show
much more dispersed ratios than healthy ones (for example, the standard
deviation of WC\_TA is 0.96 in the bankrupt group and only 0.29 in the
non-bankrupt group). Welch's version \citep{welch1947} uses each group's
own variance and adjusts the degrees of freedom accordingly. Under the null hypothesis that the two
group means are equal, the test statistic is:
\begin{equation}
t = \frac{\bar{x}_A - \bar{x}_B}{\sqrt{s_A^2/n_A + s_B^2/n_B}}
\label{eq:welch}
\end{equation}
where $\bar{x}$, $s^2$, and $n$ denote the mean, variance, and sample size
of each group ($A$ = non-bankrupt, $B$ = bankrupt). A large $|t|$ and small
$p$-value mean the observed gap between the group means is unlikely to be
due to chance.

\textbf{Why Cohen's $d$?} The $t$-statistic grows with sample size, so
even tiny differences become significant in large samples. Cohen's $d$
\citep{cohen1988} expresses the group difference in units of the pooled
standard deviation, making it independent of sample size. Benchmarks:
$|d| \approx 0.2$ is small, $0.5$ is medium, $0.8$ is large.

Each ratio is winsorized at the 1st and 99th percentiles within each group
before computing the statistics. Without this step, a small number of extreme
outliers would distort the means and inflate the standard deviations.

\begin{table}[ht]
\centering
\caption{Altman ratio means by bankruptcy status at H1.\textsuperscript{*}}
\label{tab:descriptives}
\small
\begin{tabular}{lcccccccc}
\toprule
 & \multicolumn{2}{c}{Non-bankrupt} & \multicolumn{2}{c}{Bankrupt} & & & \\
\cmidrule(lr){2-3}\cmidrule(lr){4-5}
Ratio & Mean & SD & Mean & SD & $t$ & $p$ & Cohen's $d$ \\
\midrule
WC\_TA    & 0.239 & 0.288 & $-0.208$ & 0.956 &  9.41 & $<0.001$ &  0.632 \\
RE\_TA    & 0.028 & 0.303 & $-0.411$ & 1.191 &  7.44 & $<0.001$ &  0.506 \\
EBIT\_TA  & 0.078 & 0.143 & $-0.160$ & 0.470 & 10.19 & $<0.001$ &  0.684 \\
BE\_TL    & 2.864 & 5.112 &  $1.876$ & 6.895 &  2.83 &    0.005 &  0.163 \\
sales\_TA & 1.525 & 0.922 &  $1.802$ & 1.735 & $-3.20$ &  0.001 & $-0.200$ \\
\bottomrule
\end{tabular}
\caption*{\footnotesize Group sizes (after dropping NaN, before winsorization):
non-bankrupt $n=5{,}498$, bankrupt $n=409$.}
\caption*{\footnotesize \textsuperscript{*} Each ratio is winsorized at the 1st and 99th percentiles within its group. Welch's $t$-test allows for unequal group variances. Cohen's $d$ is the standardized mean difference ($|d|>0.2$ small, $>0.5$ medium, $>0.8$ large). Source: author's calculations (code in Annex 1).}
\end{table}

Three patterns emerge from the table.

\paragraph{Four of five ratios move in the expected direction.} For WC\_TA,
RE\_TA, EBIT\_TA, and BE\_TL, the mean is higher among non-bankrupt firms
than among bankrupt firms, as theory predicts: bankrupt firms hold less
working capital relative to assets, have accumulated losses rather than
retained earnings, generate negative operating profits, and have lower equity
relative to their liabilities. All four differences are statistically
significant at the 1\% level.

\paragraph{The strongest separators are the profitability ratios.} EBIT\_TA
has the largest Cohen's $d$ at 0.68 (medium-to-large), followed by WC\_TA
at 0.63 and RE\_TA at 0.51. This matches the finding in the original
\cite{altman1968} paper that operating and cumulative profitability are the
dominant predictors. BE\_TL has a $d$ of only 0.16 (small), which is
consistent with its high variability, the standard deviations (5.1 and
6.9) are much larger than the means (2.9 and 1.9), so the within-group
spread swamps the across-group gap.

\paragraph{The sales-to-assets ratio moves in the wrong direction.} In
Altman's original Z-score, $X_5$ (sales over total assets) carries a
positive coefficient: high asset turnover is taken as a sign of efficient
operations and therefore lower bankruptcy risk. In our Polish H1 sample, the
mean of sales\_TA is actually \emph{higher} among bankrupt firms (1.80 vs
1.53, $t = -3.20$, $p = 0.001$). The most likely explanation is that distressed firms reduce their asset
base in the final year before filing, selling or writing down assets, which
mechanically raises sales-to-assets even as the firm deteriorates. This is
a real contextual difference between the original U.S.\ 1960s setting and
the Polish data, and the re-estimated Altman model is expected to produce a
different sign on this variable.

Overall, the descriptive analysis confirms that the Altman ratios carry real
information about bankruptcy in the Polish data. Four of five ratios show
medium-to-large effects in the expected direction, providing a solid
foundation for the re-estimated models. The behavior of sales\_TA sets up
a direct test of whether re-estimation adapts the framework to a new
institutional context.

\clearpage
% =============================================================================
\section{Methodology}\label{sec:methodology}
% =============================================================================

This section sets out the models we compare and the single evaluation
protocol we apply to all of them. The central design principle is fairness:
every model is trained and tested under identical conditions, so that any
performance difference reflects the model itself and not differences in
sampling, feature sets, or evaluation rules.

\subsection{Notation and the prediction task}

Let $i$ index a firm-year observation. The outcome is binary, $y_i = 1$ if
the firm goes bankrupt within the relevant horizon and $y_i = 0$ otherwise.
Each observation has a feature vector $\mathbf{x}_i$ containing the 64
financial ratios plus the 11 missingness indicators introduced in
Section~\ref{sec:missing}, for 75 candidate predictors in total. Every model produces an estimated bankruptcy probability
$\hat{p}_i \in [0, 1]$, and we evaluate these under a common protocol.
Each of the five horizons is treated as a separate prediction problem,
with H1 as the primary horizon.

\subsection{Model specifications}

We compare six models in three families of increasing flexibility. The
aim of this subsection is to clarify what each model is optimizing.

\paragraph{Re-estimated econometric benchmarks.} The first family consists of
the three classical models from Section~\ref{sec:lit}, re-estimated on the
Polish data. We use \emph{re-estimated} coefficients rather than the original published
weights because those weights were calibrated on U.S.\ firms in the 1960s
and do not transfer directly to Polish firms from the 2000s.

\begin{itemize}
\item \textbf{Altman (logit).} A logistic regression on the five Altman ratios
(WC\_TA, RE\_TA, EBIT\_TA, BE\_TL, sales\_TA). We deliberately use a logit
rather than Altman's original discriminant analysis: the logit makes no
normality assumption and outputs a probability on the same $[0,1]$ scale as
every other model.

\item \textbf{Ohlson (logit).} A logistic regression on the seven recoverable
O-score variables (log\_TA, TL\_TA, WC\_TA, CA\_STL, net\_profit\_TA,
NP\_DA\_TL, equity\_TA) plus the OENEG indicator. The two unavailable
variables, INTWO and CHIN, require firm-level lags that the dataset does not
permit.

\item \textbf{Zmijewski (probit).} A probit regression on the three X-score
variables (net\_profit\_TA, TL\_TA, CA\_STL). We retain the probit link for
the coefficient table to stay faithful to the original specification; in the
cross-validated scoring harness, the same three-variable specification is also
run as a logit to keep probability extraction and calibration identical.
\end{itemize}

\paragraph{Regularized linear models.} The second family bridges the
econometric and ML approaches: linear models that use all 75 features but
control the resulting overfitting and multicollinearity through a penalty on
the coefficients. We include two, which differ in the kind of penalty they
apply.

\begin{itemize}
\item \textbf{LASSO logit.} A logistic regression on all 75 features with an
$L_1$ penalty. The estimator solves
\begin{equation}
\min_{\boldsymbol{\beta}} \left[-\ell(\boldsymbol{\beta}; X, y)
+ \lambda \sum_{j=1}^{p} |\beta_j|\right],
\label{eq:lasso-objective}
\end{equation}
where $\ell(\boldsymbol{\beta}; X, y)$ is the logistic log-likelihood and
$\lambda = 1/C$ is the penalty strength. When $\lambda>0$, weak coefficients
are shrunk toward zero, and many are set exactly to zero. The LASSO performs automatic variable selection while staying a linear model.
The variables that survive are reported in Section~6.

\item \textbf{Ridge logit.} A logistic regression on the same 75 features with
an $L_2$ penalty,
\begin{equation}
\min_{\boldsymbol{\beta}} \left[-\ell(\boldsymbol{\beta}; X, y)
+ \lambda \sum_{j=1}^{p} \beta_j^2\right].
\label{eq:ridge-objective}
\end{equation}
The squared penalty shrinks the coefficients toward zero but, unlike the
LASSO, does not set any of them exactly to zero. Ridge keeps all 75 features
in the model and is the standard tool for stabilizing a regression when the
predictors are numerous and highly correlated, as the financial ratios here
are. It is included for a specific reason: it is the strongest \emph{linear}
competitor that uses exactly the same information as the tree ensembles, so
comparing it to XGBoost isolates the value of nonlinearity from the value of a
broad feature set. We use $C=1.0$, the standard default; performance is stable
across the range $C \in [0.1, 10]$.
\end{itemize}

\paragraph{Machine-learning models.} The third family consists of two
tree-based ensembles that can capture nonlinearities and interactions.

\begin{itemize}
\item \textbf{Random forest.} A random forest averages many decision trees. For
firm $i$, the predicted probability is
\begin{equation}
\hat{p}_i = \frac{1}{B}\sum_{b=1}^{B} T_b(\mathbf{x}_i),
\label{eq:rf-prediction}
\end{equation}
where $B$ is the number of trees and $T_b(\mathbf{x}_i)$ is the probability
prediction from tree $b$. The trees are decorrelated by two sources of
randomness: each tree is grown on a bootstrap resample of the training data,
and each split considers only a random subset of the available features. This
reduces variance relative to a single tree while preserving the ability to
model nonlinear thresholds.

\item \textbf{Gradient-boosted trees (XGBoost).} XGBoost builds an additive
ensemble sequentially. At boosting step $t$, it adds a new tree $f_t$ by
minimizing a regularized objective of the form
\begin{equation}
\mathcal{L}^{(t)} = \sum_{i=1}^{n} \ell\!\left(y_i,
\hat{y}_i^{(t-1)} + f_t(\mathbf{x}_i)\right) + \Omega(f_t),
\qquad
\Omega(f_t)=\gamma T + \frac{1}{2}\lambda \sum_{j=1}^{T} w_j^2,
\label{eq:xgb-objective}
\end{equation}
where $T$ is the number of leaves in the new tree and $w_j$ are the leaf
weights \citep{chen2016}. The new tree is fitted to the gradient of the loss
with respect to the current ensemble prediction, while the penalty discourages
overly complex trees. This makes XGBoost a regularized boosting method rather
than just a sequence of unconstrained trees.
\end{itemize}

All models within a given comparison use the same fold assignments and the
same preprocessing rules. When feature sets differ, as in the classical versus
full-feature comparisons, we interpret the gap explicitly as a mixture of
feature-breadth and algorithmic effects rather than attributing everything to
the algorithm.

\subsection{How the tree-based models work}
\label{sec:ml-theory}

The formulas above define what each model optimizes, but the two tree-based
methods deserve a short conceptual explanation, since they are the models that
drive most of the results and are less familiar than logistic regression.

\paragraph{Decision trees.} A decision tree is the building block of both
ensemble methods. It splits the sample step by step: at each node it looks for
the single feature and threshold that best separate bankrupt from healthy
firms, for example ``operating expenses to liabilities above some cutoff,''
then repeats the procedure inside each resulting group. The quality of a split
is measured by how much it reduces an impurity criterion such as the Gini
index or cross-entropy, both of which are smallest when a node contains firms
of only one class. The appeal of a tree is that it captures nonlinearities and
interactions automatically: a ratio can matter only in some region of its
range, or only in combination with another ratio, without the modeler
specifying that in advance. The drawback is that a single deep tree is
unstable: small changes in the training data can produce a very different
tree, which means high variance and a tendency to overfit. The two ensemble
methods we use are different ways of controlling that variance.

\paragraph{Bagging and the random forest.} A random forest grows many trees
and averages them, as in Equation~\eqref{eq:rf-prediction}. Two devices make
the averaging effective \citep{breiman2001}. First, each tree is trained on a
bootstrap resample of the data, so the trees see slightly different samples.
Second, at each split only a random subset of the features is considered, which
stops a few strong predictors from dominating every tree and makes the trees
less correlated with one another. Because averaging reduces variance most when
the averaged items are weakly correlated, these two sources of randomness are
what let the forest keep the flexibility of trees while removing much of their
instability.

\paragraph{Gradient boosting and XGBoost.} Boosting takes the opposite route to
bagging. Instead of building many independent trees in parallel and averaging,
it builds trees in sequence, each one correcting the errors left by the
ensemble so far \citep{friedman2001}. At each step a new small tree is fitted
to the gradient of the loss with respect to the current predictions, so it
focuses on the firms the model currently gets wrong. XGBoost is a modern,
regularized implementation of this idea \citep{chen2016}: as
Equation~\eqref{eq:xgb-objective} shows, it adds an explicit penalty on the
number of leaves and the size of the leaf weights, and it uses a second-order
(Newton-style) approximation of the loss when choosing splits. The penalty is
what keeps a long sequence of trees from overfitting, and it is the main reason
boosting can be pushed to high accuracy on this kind of data while remaining
reasonably well calibrated. We interpret these models after the fact using SHAP
values in Section~\ref{sec:shap}, which is the standard tool for explaining
tree-ensemble predictions \citep{lundberg2017}.

\subsection{The unified evaluation protocol}

\paragraph{Stratified five-fold cross-validation.} We randomly partition each
horizon's data into five folds, stratified on the bankruptcy label so that
each fold contains approximately the same default rate. Each fold serves once
as the validation set while the other four are used for training; performance
is averaged across the five folds. Stratification is essential here because
the small number of bankrupt firms (271-515 per horizon) means a purely
random split could leave some folds with very few positives.

\paragraph{Class weighting instead of resampling.} To handle the class
imbalance, we weight the minority (bankrupt) class in the loss function by
the inverse of its frequency, rather than resampling the data. Concretely,
each class $c$ receives weight $w_c = N / (2 N_c)$, where $N$ is the total
number of observations and $N_c$ the number in class $c$, so that the two
classes contribute equally to the loss. This uses all of the data and, unlike
undersampling or SMOTE, does not distort the base rate as severely, keeping
the predicted probabilities closer to correctly calibrated
\citep{king2001}. We test resampling alternatives only as a robustness check
in Section~7.

\paragraph{Leakage-free imputation.} Missing ratio values are filled with the
median, but the median is computed on the \emph{training portion of each fold
only} and then applied to the validation fold. Computing the median on the
full dataset before splitting would let information from the validation fold
leak into training, a common but subtle error. The 11 missingness indicators
are derived directly from the raw missing-value pattern, involve no
estimation, and are therefore safe to construct once for the whole dataset.

\paragraph{Identical metrics.} Every model is scored on the full battery of
metrics described next, computed identically.

\subsection{Evaluation metrics}

Because no single number captures all the properties of a probabilistic
classifier, we report metrics covering four distinct dimensions:
discrimination, calibration, operating-point behavior, and economic cost.

\paragraph{Discrimination.} These metrics measure how well a model
\emph{ranks} firms by risk, regardless of the absolute level of the predicted
probabilities.
\begin{itemize}
\item \textbf{ROC-AUC.} The area under the receiver operating characteristic
curve. It equals the probability that a randomly chosen bankrupt firm
receives a higher predicted probability than a randomly chosen healthy firm.
A value of 0.5 is random guessing and 1.0 is perfect ranking.

\item \textbf{PR-AUC.} The area under the precision-recall curve, also called
average precision. \emph{Precision} is the share of firms flagged as bankrupt
that truly are bankrupt; \emph{recall} is the share of true bankruptcies that
the model flags. As \cite{saito2015} show, PR-AUC is more informative than
ROC-AUC under heavy class imbalance, because the ROC curve can look strong
even when most positive predictions are false alarms. The no-skill baseline
for PR-AUC equals the base rate (around 0.04-0.07 here), so values well above
that indicate genuine signal.

\item \textbf{KS statistic.} The Kolmogorov-Smirnov statistic is the largest
vertical gap between the cumulative score distributions of the bankrupt and
healthy firms, equivalently $\max_t \big(\text{TPR}(t) - \text{FPR}(t)\big)$
over thresholds $t$. It is a standard summary of separation in credit scoring.
\end{itemize}

\paragraph{Calibration.} These metrics measure whether the predicted
probabilities are correct in \emph{absolute} terms, whether firms assigned
a 10\% probability actually default about 10\% of the time.
\begin{itemize}
\item \textbf{Brier score.} The mean squared difference between predicted
probability and outcome \citep{brier1950},
\begin{equation}
\text{Brier} = \frac{1}{N} \sum_{i=1}^{N} (\hat{p}_i - y_i)^2 .
\label{eq:brier}
\end{equation}
It is a \emph{proper scoring rule}: it is minimized by reporting the true
probabilities, and it rewards both good ranking and good calibration. Lower
is better.

\item \textbf{Expected Calibration Error (ECE).} We sort predictions into
$M = 10$ equal-width probability bins and average the gap between the observed
default rate and the mean predicted probability in each bin, weighted by bin
size:
\begin{equation}
\text{ECE} = \sum_{m=1}^{M} \frac{|B_m|}{N}
\,\big|\, \text{acc}(B_m) - \text{conf}(B_m) \,\big| ,
\label{eq:ece}
\end{equation}
where $|B_m|$ is the number of observations in bin $m$, $\text{acc}(B_m)$ is
the observed default rate in the bin, and $\text{conf}(B_m)$ is the mean
predicted probability in the bin \citep{guo2017}. ECE $= 0$ means perfectly
calibrated probabilities.

\item \textbf{Reliability diagram.} A plot of observed default rate against
mean predicted probability across bins; a perfectly calibrated model lies on
the 45-degree line. These are reported in Section~5.
\end{itemize}

\paragraph{Operating-point metrics.} To turn probabilities into a yes/no
bankruptcy call we need a threshold. We choose the threshold that maximizes
the F1 score on the bankrupt class, the harmonic mean of precision and
recall, and at that threshold we report precision, recall (sensitivity),
the false-positive rate (Type-I error), and the false-negative rate (Type-II
error). Reporting both error types separately matters because they have very
different economic consequences.

\paragraph{Economic cost.} Discrimination and calibration metrics treat the
two error types symmetrically, but in credit risk they are not. Following
\cite{altman2007}, we assume a false negative (missing a firm that goes on to
fail) is 35 times as costly as a false positive (a false alarm). We sweep the
classification threshold to find the one that minimizes the total cost
\begin{equation}
C = c_{\text{FN}} \cdot \text{FN} + c_{\text{FP}} \cdot \text{FP},
\qquad c_{\text{FN}} : c_{\text{FP}} = 35 : 1,
\label{eq:cost}
\end{equation}
and report this minimized cost normalized by the cost of the trivial rule that
predicts every firm to be healthy ($c_{\text{FN}} \times$ the number of true
bankruptcies). A normalized cost below one means the model improves on that
naive baseline; the 35:1 ratio is varied in robustness checks.

\paragraph{Post-hoc recalibration.} Tree ensembles in particular are known to
produce poorly calibrated probabilities even when they rank well
\citep{niculescu2005}. We therefore also report calibration after a post-hoc
fix: within each training fold we hold out 25\% of the data, fit an
\emph{isotonic regression}, a non-parametric, monotonic mapping from raw
scores to calibrated probabilities, on that held-out slice, and apply the
mapping to the validation fold. Because the recalibration is fit on data not
used to train the model and not part of the validation fold, it introduces no
leakage. Comparing calibration before and after this step isolates how much
of any calibration problem is fixable after the fact.

\subsection{Implementation}

All models are implemented in Python using scikit-learn (logistic regression,
random forest, isotonic regression, cross-validation) and XGBoost
(gradient-boosted trees); the probit benchmark uses statsmodels. The
cross-validation harness fits all preprocessing inside each fold, supports
optional feature scaling for the linear models, and computes the full metric
battery in a single reusable function applied identically to every model. A
fixed random seed is used throughout so that the fold assignments, and hence
the reported results, are reproducible.

\clearpage
% =============================================================================
\section{Results I: Econometric Baselines}\label{sec:results-econ}
% =============================================================================

We begin with the three classical models, re-estimated on the Polish data at
the one-year horizon (H1). For each model we report two things: the estimated
coefficients with their statistical significance (which tells us \emph{which}
ratios matter and in \emph{which} direction), and the cross-validated
predictive performance under the unified protocol (which tells us how well the
model actually forecasts). These results serve as the benchmark that the
machine-learning models in the next section must beat.

\subsection{Estimated coefficients}

Table~\ref{tab:econ-coef} reports the coefficients. To make the magnitudes
comparable within each model, every continuous predictor was standardized to
have mean zero and unit standard deviation before estimation, so each
coefficient measures the effect of a one-standard-deviation increase in that
ratio. For the two logit models we also report the odds ratio
($\exp(\text{coef})$), which is the multiplicative effect on the odds of
bankruptcy: a value below one means the ratio reduces default risk, and above
one means it raises it.

\begin{table}[ht]
\centering
\caption{Re-estimated econometric models at horizon H1.\textsuperscript{*}}
\label{tab:econ-coef}
\small
\begin{tabular}{lrrrr}
\toprule
Variable & Coef. & Std.\ err. & Odds ratio & Sig. \\
\midrule
\multicolumn{5}{l}{\textit{Panel A: Altman (logit, 5 ratios). Pseudo-$R^2 = 0.156$}} \\
\midrule
Constant            & $-2.977$ & 0.065 & 0.051 & *** \\
WC\_TA              & $-0.358$ & 0.057 & 0.699 & *** \\
RE\_TA              & $-0.029$ & 0.051 & 0.971 &     \\
EBIT\_TA            & $-0.713$ & 0.058 & 0.490 & *** \\
BE\_TL              & $ 0.054$ & 0.066 & 1.056 &     \\
sales\_TA           & $ 0.129$ & 0.049 & 1.137 & *** \\
\midrule
\multicolumn{5}{l}{\textit{Panel B: Ohlson (logit, 7 ratios + OENEG). Pseudo-$R^2 = 0.183$}} \\
\midrule
Constant            & $-3.047$ & 0.073 & 0.047 & *** \\
log\_TA             & $-0.452$ & 0.060 & 0.637 & *** \\
TL\_TA              & $-0.252$ & 0.350 & 0.777 &     \\
WC\_TA              & $-0.462$ & 0.081 & 0.630 & *** \\
CA\_STL             & $ 0.220$ & 0.073 & 1.246 & *** \\
net\_profit\_TA     & $-0.483$ & 0.065 & 0.617 & *** \\
NP\_DA\_TL          & $-0.366$ & 0.131 & 0.694 & *** \\
equity\_TA          & $-0.395$ & 0.344 & 0.674 &     \\
OENEG               & $-0.701$ & 0.277 & 0.496 & **  \\
\midrule
\multicolumn{5}{l}{\textit{Panel C: Zmijewski (probit, 3 ratios). Pseudo-$R^2 = 0.154$}} \\
\midrule
Constant            & $-1.644$ & 0.029 &,   & *** \\
net\_profit\_TA     & $-0.346$ & 0.027 &,   & *** \\
TL\_TA              & $ 0.217$ & 0.027 &,   & *** \\
CA\_STL             & $ 0.056$ & 0.030 &,   & *   \\
\bottomrule
\end{tabular}
\caption*{\footnotesize \textsuperscript{*} Continuous predictors are standardized, so each
coefficient is the effect of a one-standard-deviation increase. Odds ratios are
reported for the two logit models. Significance: \textsuperscript{***}~$p<0.01$,
\textsuperscript{**}~$p<0.05$, \textsuperscript{*}~$p<0.10$. Source: author's
calculations.}
\end{table}

Three points are worth drawing out.

\paragraph{Profitability is the dominant predictor in every model.} The
operating-profitability ratio EBIT\_TA in Altman, and the net-profitability
ratio net\_profit\_TA in Ohlson and Zmijewski, all carry large, negative,
highly significant coefficients. In the Altman model, a one-standard-deviation
increase in EBIT\_TA roughly halves the odds of bankruptcy (odds ratio 0.49).
This matches both the descriptive statistics of Section~\ref{sec:descriptives}
and the broader literature, in which profitability consistently ranks among
the strongest distress signals.

\paragraph{The sales-to-assets sign flip is confirmed.} As anticipated in
Section~\ref{sec:descriptives}, the coefficient on sales\_TA in the Altman
model is positive and significant (odds ratio 1.14), the opposite of the
positive-for-health role it plays in Altman's original 1968 Z-score. In the
Polish data, higher asset turnover is associated with \emph{higher}, not
lower, default risk. The most plausible explanation is that firms approaching
failure shrink their asset base, through asset sales or write-downs, in
their final reporting year, which mechanically raises sales relative to
assets. This is a clear example of a classical coefficient that does not
transfer unchanged to a new institutional and temporal setting, and it
illustrates why we re-estimate the models rather than applying the original
weights.

\paragraph{Some ratios lose significance once others are included.} RE\_TA and
BE\_TL are individually informative in the descriptive analysis but become
insignificant in the multivariate Altman model, because their information
overlaps with the stronger profitability and liquidity ratios. In the Ohlson model, CA\_STL enters with a positive sign, which runs against
simple intuition; sign reversals like this are common when several correlated
ratios are included together. We treat the individual Ohlson coefficients with
caution. What matters for our purposes is predictive performance.

\subsection{Predictive performance}

Table~\ref{tab:econ-cv} reports the cross-validated performance of the three
models under the unified protocol: stratified five-fold cross-validation,
class-weighted loss, and leakage-free imputation. All numbers are means across
the five folds.

\begin{table}[ht]
\centering
\caption{Performance of the econometric baselines at horizon H1.\textsuperscript{*}}
\label{tab:econ-cv}
\small
\begin{tabular}{lrrrrrrr}
\toprule
Model & AUC & PR-AUC & KS & Brier & Recall & Precision & Norm.\ cost \\
\midrule
Altman (5)        & 0.772 & 0.265 & 0.478 & 0.192 & 0.476 & 0.294 & 0.335 \\
Ohlson (7+OENEG)  & 0.795 & 0.297 & 0.507 & 0.176 & 0.463 & 0.324 & 0.312 \\
Zmijewski (3)     & 0.776 & 0.286 & 0.460 & 0.189 & 0.478 & 0.303 & 0.322 \\
\bottomrule
\end{tabular}
\caption*{\footnotesize \textsuperscript{*} Means across five folds. AUC and PR-AUC measure ranking ability; Brier and ECE measure calibration (lower is better); recall and precision are at the F1-optimal threshold; normalized cost is relative to the naive ``predict everyone healthy'' rule (below one is better). Base rate 6.94\%. Source: author's calculations (code in Annex 1).}
\end{table}

The three models perform similarly, with Ohlson modestly ahead on most
metrics, an AUC of 0.795 against 0.772 for Altman and 0.776 for Zmijewski,
and the best PR-AUC, Brier score, and normalized cost. This ordering is
sensible: Ohlson has the most predictors (eight versus five and three) and the
highest in-sample fit (pseudo-$R^2$ of 0.18). The differences between the three are small relative to fold-to-fold
variation, so we do not read too much into the ranking within this family.

Two broader observations set up the comparison with machine learning. First,
all three models clearly beat random guessing: an AUC near 0.78-0.80 and a
PR-AUC of roughly 0.27-0.30 against a base rate of 0.07 represent real
predictive signal. Second, all three are poorly calibrated out of the box, with
expected calibration errors above 0.30 (not shown in the table but discussed
in Section~\ref{sec:methodology}); this is a direct consequence of training
with class weighting, which improves ranking and recall at the cost of
inflating the predicted probabilities. We return to this calibration issue
when we evaluate the machine-learning models, where post-hoc recalibration
plays an important role. The headline benchmark to carry forward is an AUC of
roughly 0.80 and a PR-AUC of roughly 0.30, achieved by the best econometric
model (Ohlson) at the one-year horizon.

Before turning to the machine-learning models, it is worth being precise about
what these econometric results measure. The three classical models use between
three and eight variables, all accounting ratios from a single filing date.
Modern econometric default models, from \cite{shumway2001} onwards, augment
these with market-based controls such as firm size, past equity returns,
idiosyncratic volatility, and Merton-style distance-to-default, which
\cite{campbell2008} show to be among the most informative predictors
available. Industry effects, as in \cite{chava2004}, capture the fact that
baseline default rates and the marginal effect of leverage differ across
sectors. None of these controls are available in the Polish dataset, because
the firms are private and anonymized. The three re-estimated models are
therefore not the strongest possible econometric specification: they are the
strongest specification feasible on this dataset. When the ML models later
outperform them, part of that gap reflects algorithmic flexibility and a
broader ratio set, but part reflects this structural disadvantage in control
variables. The feature-breadth analysis in Section~\ref{sec:feature-breadth}
is designed precisely to separate these effects.

\clearpage
% =============================================================================
\section{Results II: Machine Learning Models and Full Comparison}\label{sec:results-ml}
% =============================================================================

We now turn to the machine-learning models and compare them, under the
unified protocol, against the econometric benchmarks just established. The
comparison is reported in stages. First, we look at the two regularized linear
models, LASSO and Ridge, both fit on all 75 features. These sit between the
hand-picked econometric models and the fully flexible tree ensembles: they use
the broad feature set but remain linear. Then we report the random forest and
XGBoost results and assemble the full seven-model comparison.

\subsection{Regularized linear models: LASSO and Ridge}

The LASSO logit is fit on all 75 features, with the $L_1$ penalty driving
weak coefficients to exactly zero. The penalty strength $C$ is fixed at 0.1;
performance is essentially flat across the range $C \in [0.05, 0.2]$,
so a single value is sufficient. The
variables that survive with non-zero coefficients are reported in
Table~\ref{tab:lasso-survivors}.

\begin{table}[ht]
\centering
\caption{Top variables selected by the LASSO logit, H1.\textsuperscript{*}}
\label{tab:lasso-survivors}
\small
\begin{tabular}{lrl}
\toprule
Variable & Coef. & Interpretation \\
\midrule
miss\_sales\_growth        & $+4.136$ & raises risk \\
miss\_OP\_finexp           & $+2.981$ & raises risk \\
miss\_TL\_op\_DA\_monthly  & $-2.051$ & lowers risk \\
sales\_STL                 & $-1.812$ & lowers risk \\
OpEx\_STL                  & $+1.708$ & raises risk \\
miss\_GP3yr\_TA            & $-1.313$ & lowers risk \\
EBITDA\_TA                 & $+0.615$ & raises risk \\
net\_profit\_TA            & $-0.581$ & lowers risk \\
BE\_TL                     & $-0.550$ & lowers risk \\
profit\_sales\_TA          & $-0.526$ & lowers risk \\
GP\_extra\_fin\_TA         & $+0.457$ & raises risk \\
CA\_inv\_STL               & $-0.442$ & lowers risk \\
equity\_TA                 & $+0.437$ & raises risk \\
CA\_TL                     & $+0.437$ & raises risk \\
const\_cap\_TA             & $-0.427$ & lowers risk \\
\bottomrule
\end{tabular}
\caption*{\footnotesize \textsuperscript{*} Penalty $C = 0.1$. Coefficients are on
standardized predictors; positive values raise the modeled bankruptcy risk,
negative values lower it. Of the 75 candidate features, 63 survived; the 15
with the largest absolute coefficient are shown. Of the 11 hand-picked
variables used by Altman, Ohlson, and Zmijewski, 8 survive the LASSO:
net\_profit\_TA, BE\_TL, equity\_TA, log\_TA, WC\_TA, TL\_TA, CA\_STL, RE\_TA.
Source: author's calculations (code in Annex 1).}
\end{table}

The most striking feature of Table~\ref{tab:lasso-survivors} is that the
\emph{three largest coefficients in the entire model are missingness
indicators rather than financial ratios}. The fact that a firm fails to
report \texttt{sales\_growth} carries a standardized coefficient of $+4.14$,
more than six times larger than the strongest financial ratio coefficient.
This is striking confirmation of the informative-missingness finding
flagged descriptively in Section~\ref{sec:missing}: \emph{whether} a
ratio is reported is a more powerful bankruptcy signal than the
ratio's value itself. The economic interpretation is that distressed firms
selectively fail to disclose certain items, presumably the most damaging
ones, and the pattern of non-disclosure is what carries the predictive
weight. Among the ratios proper, the LASSO confirms the central role of
profitability variants (\texttt{EBITDA\_TA}, \texttt{net\_profit\_TA},
\texttt{profit\_sales\_TA}) and pulls in a small number of additional
liquidity and capital-structure measures. Of the 11 variables used in the
canonical econometric models, 8 survive the penalty, with the remaining 3
absorbed by overlapping predictors.

The cross-validated predictive performance of the LASSO is reported in the
combined table below; here we note only that the LASSO already represents a
substantial jump over the econometric benchmarks, with an AUC of 0.895 and a
PR-AUC of 0.514, roughly a 10-point gain in AUC and an almost 22-point gain
in PR-AUC over Ohlson. This jump is attributable to two things: access to the
full 75-feature set, and explicit use of the missingness indicators that the
hand-picked benchmarks cannot use.

The Ridge logit is fit on the same 75 features but with an $L_2$ penalty,
which shrinks the coefficients without setting any to zero, so all 75 features
stay in the model. We use the standard default penalty ($C = 1.0$); as noted
in Section~\ref{sec:methodology}, the result is stable across a wide range of
$C$. Ridge is included as the strongest purely linear competitor that sees
exactly the same information as the tree ensembles. Its performance is very
close to the LASSO's, with an AUC of 0.877 and a PR-AUC of 0.448. The small
gap between the two reflects only the different penalty: the LASSO discards
the weakest features, while Ridge keeps them but shrinks them, and on this
dataset the two strategies land in almost the same place. The point of
including Ridge is what comes next: even a linear model that uses the entire
feature set, with multicollinearity controlled by the penalty, still falls
well short of the tree ensembles. That gap cannot be explained by feature
count, since Ridge and XGBoost use the same 75 features, so it must come from
nonlinearity. We return to this comparison in Section~\ref{sec:feature-breadth}.

\subsection{Random forest and XGBoost}

The two tree-based ensembles use the same 75-feature set as the LASSO. The
random forest is fit with 300 trees and the default $\sqrt{p}$ features
sampled at each split, with no depth cap and no hyperparameter tuning.
XGBoost uses moderately conservative settings (400 trees, learning rate
0.05, maximum depth 4, subsampling 0.8 on both observations and features,
$L_2$ regularization of 1.0), with no extensive hyperparameter search , 
the goal of the comparison is to report a defensible default rather than to
chase the maximum achievable AUC. Both models use built-in class-imbalance
handling consistent with the rest of the project: the random forest uses
inverse-frequency class weights, while XGBoost uses its native
\texttt{scale\_pos\_weight} parameter set to the ratio of healthy to
bankrupt firms.

\subsection{Full seven-model comparison}

Table~\ref{tab:full-comparison} reports the cross-validated performance of
all seven models at H1, both before and after post-hoc isotonic
recalibration. The recalibration step fits an isotonic regression on a 25\%
holdout slice of each training fold and applies the resulting monotonic
mapping to the predicted probabilities of that fold's validation set; this
isolates how much of any calibration problem is fixable after the fact.

\begin{table}[ht]
\centering
\caption{Performance of all seven models at horizon H1, before and after
isotonic recalibration.\textsuperscript{*}}
\label{tab:full-comparison}
\small
\begin{tabular}{lrrrrrrr}
\toprule
Model & AUC & PR-AUC & KS & Brier & ECE & Recall & Norm.\ cost \\
\midrule
\multicolumn{8}{l}{\textit{Panel A: raw probabilities (class-weighted training, no post-hoc recalibration)}} \\
\midrule
Altman (5)         & 0.772 & 0.265 & 0.478 & 0.192 & 0.359 & 0.476 & 0.335 \\
Ohlson (7+OENEG)   & 0.795 & 0.297 & 0.507 & 0.176 & 0.327 & 0.463 & 0.312 \\
Zmijewski (3)      & 0.776 & 0.286 & 0.460 & 0.189 & 0.356 & 0.478 & 0.322 \\
LASSO logit (75)   & 0.895 & 0.514 & 0.702 & 0.104 & 0.205 & 0.622 & 0.181 \\
Ridge logit (75)   & 0.877 & 0.448 & 0.673 & 0.111 & 0.205 & 0.639 & 0.209 \\
Random Forest      & 0.938 & 0.627 & 0.748 & 0.040 & 0.021 & 0.646 & 0.142 \\
XGBoost            & 0.960 & 0.810 & 0.792 & 0.030 & 0.023 & 0.683 & 0.124 \\
\midrule
\multicolumn{8}{l}{\textit{Panel B: after post-hoc isotonic recalibration}} \\
\midrule
Altman (5)         & 0.765 & 0.237 & 0.460 & 0.059 & 0.013 & 0.522 & 0.350 \\
Ohlson (7+OENEG)   & 0.791 & 0.253 & 0.488 & 0.058 & 0.013 & 0.541 & 0.326 \\
Zmijewski (3)      & 0.763 & 0.244 & 0.438 & 0.059 & 0.016 & 0.515 & 0.351 \\
LASSO logit (75)   & 0.876 & 0.457 & 0.644 & 0.046 & 0.011 & 0.615 & 0.231 \\
Ridge logit (75)   & 0.872 & 0.423 & 0.658 & 0.047 & 0.012 & 0.576 & 0.221 \\
Random Forest      & 0.928 & 0.557 & 0.723 & 0.041 & 0.017 & 0.583 & 0.161 \\
XGBoost            & 0.951 & 0.747 & 0.767 & 0.028 & 0.012 & 0.624 & 0.140 \\
\bottomrule
\end{tabular}
\caption*{\footnotesize \textsuperscript{*}~Means across five folds. Recall and precision are
measured at the F1-optimal threshold; normalized cost is relative to the naive
``predict everyone healthy'' rule under a 35:1 false-negative to false-positive
cost ratio. The two regularized linear models (LASSO and Ridge) use all 75
features. Source: author's calculations (code in Annex 1).}
\end{table}

The full table now lets us read the comparison as a clear progression from
simple to complex. The three classical models sit at AUC 0.77-0.80. The two
regularized linear models, which use the full 75-feature set, jump to AUC
0.88-0.90: a large gain that comes purely from giving a linear model more
information. The two tree ensembles go further still, to 0.94-0.96, by adding
nonlinearity on top of the broad feature set. Ridge and LASSO land very close
to each other (AUC 0.877 versus 0.895), which is expected since they differ
only in the type of penalty; the important point is that both clearly beat the
classical models and both clearly lose to the tree ensembles. Three results
are worth drawing out in more detail.

\paragraph{The machine-learning advantage is large, not marginal.} The gap
between XGBoost (AUC 0.960, PR-AUC 0.810) and Ohlson (AUC 0.795,
PR-AUC 0.297) is 17 AUC points and over 50 PR-AUC points. This is far
larger than the 5-10 point gap typically reported in ML-vs-econometric
comparison studies, and it survives the unification of the evaluation
protocol, so it cannot be dismissed as an artifact of differential class
balancing or evaluation rules. The PR-AUC gap is particularly striking: under
the 7\% base rate, the no-skill PR-AUC is 0.07, so XGBoost's 0.810 represents
a nearly twelvefold lift over chance, while Ohlson's 0.297 is only a
fourfold lift. In the rare-event regime that bankruptcy prediction is, this
is the more economically meaningful measure than ROC-AUC.

\paragraph{The advantage is concentrated where it matters: rare-event
discrimination, not just any-event discrimination.} The ROC-AUC gap of 17
points is large; the PR-AUC gap of 51 points is much larger. The
interpretation is that the tree ensembles do not just slightly improve
overall ranking, they substantially reduce false alarms among healthy
firms while still catching a large share of bankrupt firms. Concretely, at
the F1-optimal threshold XGBoost achieves precision of 0.83 and recall of
0.68 on the bankrupt class, while Ohlson achieves precision of 0.32 and
recall of 0.46. For every ten firms XGBoost flags, eight are truly distressed;
for every ten Ohlson flags, only three are.

\paragraph{Calibration and discrimination separate cleanly under
recalibration.} Panel B of Table~\ref{tab:full-comparison} confirms the
methodological prediction made in Section~\ref{sec:methodology}: post-hoc
isotonic recalibration dramatically improves the Brier score and the
expected calibration error of the class-weighted linear models (the ECE for
the econometric baselines drops from above 0.32 to around 0.013), while
leaving their discrimination metrics almost unchanged. The tree ensembles,
by contrast, are already well-calibrated out of the box, Random Forest's
ECE is 0.021 even before any post-hoc fix, and XGBoost's is 0.023, and
recalibration leaves them essentially unchanged. This separates the two
problems cleanly: the class-weighted linear models are systematically
overconfident about absolute default probabilities but rank firms well;
the tree ensembles do both well at the same time. The relative ordering of the six models is the same before and after
recalibration: XGBoost is best on every metric, trees beat the LASSO which
beats the econometric benchmarks, so none of the conclusions depend on
calibration artifacts.

\paragraph{Some headline numbers tell the story.} At H1, on the unified
protocol, our best econometric model (Ohlson) flags bankrupt firms at the
F1-optimal threshold with 32\% precision and 46\% recall; our best ML model
(XGBoost) flags them with 83\% precision and 68\% recall. In economic-cost
terms, under a 35:1 false-negative-to-false-positive cost ratio, XGBoost
achieves 12\% of the cost of the trivial ``predict everyone healthy'' rule,
while Ohlson achieves 31\%. The improvement from a classical accounting
model to a modern tree ensemble is real, large, and concentrated exactly
where rare-event practitioners need it.

\subsection{Comparison with published results on the same dataset}
\label{sec:litbenchmark}

The numbers above are internally consistent within our protocol, but the natural
next question is how they sit relative to what other researchers have found on
the same Polish dataset. This matters for three reasons. First, it is a pipeline
check: if our XGBoost were far below published results, something would be wrong
with our implementation. Second, it locates the classical-model baselines in the
existing literature, showing whether our Altman and Ohlson re-estimates are
representative of how these models perform when applied outside their original
samples. Third, and most substantively, it shows what the existing literature
has and has not measured, which is where the original contribution of this paper
lies.

Table~\ref{tab:litbenchmark} collects the main published AUC results at the
one-year-ahead horizon, and Figure~\ref{fig:litbenchmark} plots them in rank
order against our key results.

\begin{table}[H]
\centering
\caption{Published results on the Polish dataset versus ours, one-year-ahead
horizon.\textsuperscript{*}}
\label{tab:litbenchmark}
\small
\begin{tabular}{llcl}
\toprule
Study & Model & AUC & Method notes \\
\midrule
\cite{zieba2016}        & Plain logistic     & 0.632 & no resampling \\
\cite{zieba2016}        & LDA                & 0.796 & linear discriminant \\
\cite{zieba2016}        & Cost-sensitive logit & 0.821 & cost-weighted \\
\cite{zieba2016}        & Random Forest      & 0.898 & 1000 trees \\
\cite{zieba2016}        & AdaCost            & 0.928 & cost-sensitive boosting \\
\cite{zieba2016}        & XGBoost            & 0.951 & native missing handling \\
\cite{zieba2016}        & EXGB (their best)  & 0.955 & boosted trees + synthetic features \\
\cite{smiti2020}        & Deep AE + Borderline-SMOTE & $>$0.95 & deep learning, oversampling \\
\midrule
\textbf{This paper}     & Ohlson logit (7)   & 0.795 & unified protocol \\
\textbf{This paper}     & Ridge logit (75)   & 0.877 & unified protocol \\
\textbf{This paper}     & Random Forest      & 0.937 & unified protocol \\
\textbf{This paper}     & XGBoost            & 0.960 & unified protocol \\
\bottomrule
\end{tabular}
\caption*{\footnotesize \textsuperscript{*} AUC is the metric reported across this
literature. Published values from \cite{zieba2016} are from its 5thYear file,
the one-year-ahead task. Differences in preprocessing mean the numbers are
broadly, not exactly, comparable. Source: author's calculations for the ``This
paper'' rows (code in Annex 1).}
\end{table}

\begin{figure}[H]
\centering
\includegraphics[width=0.92\textwidth]{fig_litbenchmark.pdf}
\caption{Our results against published benchmarks on the Polish dataset at the
one-year-ahead horizon. Grey bars are results from \cite{zieba2016}; the blue
bar is our econometric baseline (Ohlson); the red bar is our best model
(XGBoost). Our XGBoost sits at the top of the published range, and our
classical baseline sits in the same band as other classical models.}
\label{fig:litbenchmark}
\end{figure}

\paragraph{How our results compare on the headline metric.}

Our XGBoost AUC of 0.960 sits at the very top of what has been reported on this
dataset, marginally above the 0.955 of \cite{zieba2016}'s best model (EXGB, an
ensemble of XGBoost trees with evolved synthetic features) and above their plain
XGBoost at 0.951. This is not a trivial result: the EXGB model in
\cite{zieba2016} uses a purpose-built feature generation step specifically
designed to improve discrimination on this dataset, while we use a
straightforward off-the-shelf XGBoost with standard hyperparameters. The fact
that the two methods reach essentially the same ceiling suggests that the
predictive information in the dataset is largely exhausted by a well-implemented
gradient boosting model, and that further gains would require richer data rather
than more sophisticated algorithms.

Our Random Forest at 0.937 exceeds the \cite{zieba2016} Random Forest at 0.898
by about 4 points. The likely explanation is that we use missingness indicators
as explicit features, which the original paper did not, and that XGBoost's
native handling of missing values, which propagates observations in both
directions at each split, is less effective than explicitly coding missingness
as a signal for a random forest. Our Ohlson logit at 0.795 is close to the LDA
result of \cite{zieba2016} at 0.796 and well above their plain logistic
regression at 0.632. The large gap between plain logistic regression and our
class-weighted logit confirms that the choice of imbalance treatment matters as
much as the algorithm choice for classical models.

\paragraph{Why the numbers are not exactly comparable.}

Three methodological differences mean the AUC values in Table~\ref{tab:litbenchmark}
cannot be compared as if they were produced by the same protocol.

\emph{Missing-value handling.} \cite{zieba2016} exploit XGBoost's built-in
ability to route missing values during tree splits. We instead impute medians
and add explicit missingness indicator columns. Our approach makes the
missingness signal available to every model, including the logit, not just to
XGBoost. This is why our missingness-indicator approach adds value even on top
of a model that already handles missing values internally.

\emph{Resampling.} \cite{smiti2020} apply Borderline-SMOTE before
cross-validation, meaning synthetic observations generated from the full
dataset's neighbourhood structure are present in both training and validation
folds. As \cite{king2001} document, this inflates predicted probabilities and
tends to inflate AUC when the synthetic samples are informative. We apply class
weighting only, fit entirely inside the training fold, so the validation fold
always consists of real, unseen observations.

\emph{Feature engineering.} The EXGB model of \cite{zieba2016} generates
synthetic ratio combinations as extra features using an evolutionary algorithm.
This is genuinely a more complex model than plain XGBoost; the fact that our
XGBoost slightly exceeds their EXGB result without this step further supports
the conclusion that the raw features plus missingness indicators already contain
most of the predictable information.

\paragraph{What the existing literature does not measure.}

Every study in Table~\ref{tab:litbenchmark} reports ROC-AUC and, in some cases,
accuracy or sensitivity. None reports PR-AUC, calibration, or a Brier
decomposition. Under the 6.9\% base rate at H1, this matters. An AUC of 0.95
by itself does not tell a practitioner how many false alarms the model generates
per correct flag, whether the predicted probabilities can be used as such, or
how the model's discrimination changes once calibration is corrected.

Our XGBoost achieves a PR-AUC of 0.810. To put this in context: the no-skill
baseline PR-AUC at the 6.9\% base rate is 0.069. Our model's rare-event
precision, the fraction of flagged firms that are genuinely distressed, reaches
83\% at a threshold calibrated to balance false positives and negatives. The
classical models, despite sitting in the same AUC band as other classical models
in the literature, achieve PR-AUC of only 0.27-0.30 and precision of 29-32\%.
This gap, invisible in the ROC plot, is the main reason a practitioner would
prefer the ML model for any application where the cost of a false alarm matters.

The calibration picture adds a further dimension that the prior literature
ignores entirely. Before isotonic recalibration, the classical models show ECE
values of 0.33-0.36, meaning their predicted probabilities deviate from the true
default rate by roughly 33-36 percentage points on average. XGBoost's ECE is
0.023. After recalibration, all models achieve ECE below 0.02, but the Brier
decomposition confirms that the sorting-power advantage of XGBoost is permanent
and cannot be removed by post-hoc adjustments. No prior study on this dataset
separates these two components, which makes our calibration results genuinely
new information rather than a replication of existing findings.

\paragraph{Summary.}

Our results sit at or above the published ceiling on this dataset on the AUC
metric that the prior literature reports. On the metrics the prior literature
does not report, PR-AUC, calibration, and Brier decomposition, our results are
the first available benchmark, and they tell a substantially different story
about relative model performance than the headline AUC numbers alone would
suggest.

\clearpage
% =============================================================================
\section{Visualizing the Comparison}\label{sec:figures}
% =============================================================================

The numbers in Table~\ref{tab:full-comparison} are easier to read as pictures.
This section presents three figures, all built from the same cross-validated
predictions: ROC curves, precision-recall curves, and calibration
(reliability) diagrams. All use the pooled out-of-fold predictions, meaning
each firm's prediction comes from a fold in which that firm was not used for
training.

\subsection{ROC and precision-recall curves}

Figure~\ref{fig:roc} shows the ROC curves. Each curve traces the trade-off between the true-positive rate (bankrupt
firms correctly flagged) and the false-positive rate (healthy firms wrongly
flagged) as the threshold moves. Closer to the top-left corner is better;
the diagonal is random guessing. The three econometric models form a lower
band, the LASSO sits above them, and the two tree ensembles are at the top.

\begin{figure}[H]
\centering
\includegraphics[width=0.72\textwidth]{fig_roc_h1.pdf}
\caption{ROC curves at horizon H1. Curves closer to the top-left corner
indicate better discrimination. The tree ensembles (XGBoost, Random Forest)
dominate the linear models, which in turn beat the hand-picked econometric
benchmarks.}
\label{fig:roc}
\end{figure}

Figure~\ref{fig:pr} shows the precision-recall curves, which are more
informative under class imbalance. Precision is the share of flagged firms
that are truly bankrupt; recall is the share of true bankruptcies the model
flags. The dashed line is the no-skill baseline at the 6.9\% base rate.
The gap between XGBoost and the rest is far larger here than in the ROC plot:
XGBoost keeps high precision even at high recall, while the econometric models
drop toward the baseline almost immediately. This is the figure that best
captures why the ML advantage matters in practice.

\begin{figure}[H]
\centering
\includegraphics[width=0.72\textwidth]{fig_pr_h1.pdf}
\caption{Precision-recall curves at horizon H1. The dashed line is the
no-skill baseline (the 6.9\% base rate). XGBoost keeps precision high across
a wide range of recall, while the econometric models drop toward the baseline
quickly. This is the most economically meaningful view for a rare-event
problem.}
\label{fig:pr}
\end{figure}

\subsection{Reliability diagrams}

Figure~\ref{fig:reliability} shows reliability diagrams for Ohlson and
XGBoost, before and after isotonic recalibration. Each diagram plots the
observed default rate against the mean predicted probability across ten
probability bins; a perfectly calibrated model lies on the 45-degree line.

Before recalibration (red), both models sit well below the diagonal, meaning
they systematically over-predict the probability of bankruptcy, an
expected consequence of the class-weighted training described in
Section~\ref{sec:methodology}. After recalibration (blue), both curves move
close to the diagonal, confirming that the miscalibration is almost entirely
fixable after the fact without changing the model's ranking. The remaining
wobble in the high-probability bins reflects the small number of firms that
receive such high scores, not a systematic calibration failure.

\begin{figure}[H]
\centering
\includegraphics[width=0.95\textwidth]{fig_reliability_h1.pdf}
\caption{Reliability diagrams at horizon H1, before and after isotonic
recalibration. Points on the dashed 45-degree line are perfectly calibrated.
Both models over-predict risk before recalibration (red, below the line) and
are brought close to calibration afterward (blue). Recalibration changes the
probabilities without changing the ranking, so discrimination metrics are
essentially unaffected.}
\label{fig:reliability}
\end{figure}

\subsection{Decomposing the Brier score}

The reliability diagrams show calibration visually. The Brier score is a
single number, but it combines two things a model can do well or badly. The decomposition of \cite{murphy1973} separates them:
\begin{equation}
\text{Brier} = \underbrace{\text{Reliability}}_{\text{calibration error}}
             - \underbrace{\text{Resolution}}_{\text{sorting power}}
             + \underbrace{\text{Uncertainty}}_{\text{fixed by the data}} .
\label{eq:brier-decomp}
\end{equation}
The \emph{reliability} term is the average squared gap between each bin's
predicted probability and the actual default rate in that bin: pure
calibration error, lower is better. The \emph{resolution} term measures how
far bin default rates spread away from the overall base rate: it captures
how well the model separates firms by risk, higher is better. The
\emph{uncertainty} term is $\bar{p}(1-\bar{p})$; it depends only on the
data, is the same for every model (0.065 at H1), and represents the score
of a model that simply predicts the base rate for all firms.

\begin{table}[ht]
\centering
\caption{Brier-score decomposition at horizon H1, before and after isotonic
recalibration. Reliability is calibration error (lower is better); resolution
is sorting power (higher is better); uncertainty (0.065) is the same for all
models. The three terms recombine to the Brier score via
Equation~\ref{eq:brier-decomp}.}
\label{tab:brier-decomp}
\small
\begin{tabular}{lcccc}
\toprule
Model & Reliability & Resolution & Uncertainty & Brier \\
\midrule
\multicolumn{5}{l}{\textit{Panel A: raw probabilities}} \\
\midrule
Altman (5)        & 0.134 & 0.007 & 0.065 & 0.192 \\
Ohlson (7+OENEG)  & 0.119 & 0.008 & 0.065 & 0.176 \\
Zmijewski (3)     & 0.131 & 0.008 & 0.065 & 0.188 \\
LASSO logit (75)  & 0.060 & 0.021 & 0.065 & 0.104 \\
Random Forest     & 0.002 & 0.026 & 0.065 & 0.041 \\
XGBoost           & 0.004 & 0.037 & 0.065 & 0.031 \\
\midrule
\multicolumn{5}{l}{\textit{Panel B: after isotonic recalibration}} \\
\midrule
Altman (5)        & 0.001 & 0.007 & 0.065 & 0.059 \\
Ohlson (7+OENEG)  & 0.001 & 0.007 & 0.065 & 0.058 \\
Zmijewski (3)     & 0.001 & 0.007 & 0.065 & 0.059 \\
Random Forest     & 0.000 & 0.024 & 0.065 & 0.041 \\
XGBoost           & 0.000 & 0.037 & 0.065 & 0.029 \\
\bottomrule
\end{tabular}
\caption*{\footnotesize Source: author's calculations (code in Annex 1).}
\end{table}

Table~\ref{tab:brier-decomp} makes two things clear. First, the high Brier
scores of the econometric models come almost entirely from the reliability
term: calibration error (0.12-0.13) is roughly fifteen to twenty times
larger than for the trees (around 0.003). Their resolution is also lower
(0.007-0.008 versus 0.026-0.037 for the trees), so they are worse on both
dimensions: miscalibrated and less discriminating. The trees win on both axes.

Second, Panel B shows exactly what isotonic recalibration does and does not
do. After recalibration, the reliability term collapses to essentially zero
for every model, calibration error is almost entirely removed, while the
resolution term barely moves (for example, XGBoost's resolution stays at
0.037). This confirms the central methodological point of the paper.
Calibration and discrimination are separate properties. Miscalibration is
fixable after the fact without sacrificing sorting power, but a model's
resolution, its genuine ability to tell risky firms from safe ones, is the
part that cannot be recovered by post-processing and must come from the model
itself. On that irreducible component, the gap
between the tree ensembles and the classical models is large and permanent.

\clearpage
% =============================================================================
\section{Performance Across Prediction Horizons}\label{sec:horizons}
% =============================================================================

All results so far use horizon H1. The dataset also provides H2 through H5,
with predictors recorded two to five years before the event. Longer horizons
are inherently harder: a firm five years from filing looks much more like a
healthy firm than one that is a year away. This section asks how each model
degrades as the horizon lengthens.

We re-run all seven models on each horizon under the same protocol used at H1:
stratified five-fold cross-validation, class weighting, leakage-free
imputation, identical metrics. Table~\ref{tab:horizon-grid} reports AUC and
PR-AUC across the full seven-model by five-horizon grid. Figure~\ref{fig:horizon}
shows the ROC-AUC and PR-AUC trajectories.

\begin{table}[ht]
\centering
\caption{Performance across prediction horizons.\textsuperscript{*}}
\label{tab:horizon-grid}
\small
\begin{tabular}{lccccc}
\toprule
Model & H1 & H2 & H3 & H4 & H5 \\
\midrule
\multicolumn{6}{l}{\textit{Panel A: ROC-AUC}} \\
\midrule
Altman (5)        & 0.772 & 0.700 & 0.679 & 0.611 & 0.675 \\
Ohlson (7+OENEG)  & 0.795 & 0.665 & 0.686 & 0.580 & 0.708 \\
Zmijewski (3)     & 0.776 & 0.698 & 0.693 & 0.633 & 0.657 \\
LASSO logit (75)  & 0.892 & 0.865 & 0.866 & 0.819 & 0.884 \\
Ridge logit (75)  & 0.877 & 0.857 & 0.821 & 0.799 & 0.889 \\
Random Forest     & 0.937 & 0.915 & 0.884 & 0.879 & 0.921 \\
XGBoost           & 0.960 & 0.938 & 0.935 & 0.927 & 0.959 \\
\midrule
\multicolumn{6}{l}{\textit{Panel B: PR-AUC}} \\
\midrule
Altman (5)        & 0.265 & 0.132 & 0.095 & 0.067 & 0.097 \\
Ohlson (7+OENEG)  & 0.297 & 0.110 & 0.095 & 0.063 & 0.120 \\
Zmijewski (3)     & 0.286 & 0.131 & 0.098 & 0.072 & 0.093 \\
LASSO logit (75)  & 0.500 & 0.337 & 0.259 & 0.175 & 0.303 \\
Ridge logit (75)  & 0.448 & 0.285 & 0.210 & 0.158 & 0.322 \\
Random Forest     & 0.621 & 0.585 & 0.491 & 0.515 & 0.647 \\
XGBoost           & 0.810 & 0.724 & 0.662 & 0.691 & 0.775 \\
\midrule
\textit{Base rate} & 0.069 & 0.053 & 0.047 & 0.039 & 0.039 \\
\bottomrule
\end{tabular}
\caption*{\footnotesize \textsuperscript{*} Means across five folds. H1 is one year ahead, H5 is five years ahead. Upper panel: ROC-AUC; lower panel: PR-AUC. Source: author's calculations (code in Annex 1).}
\end{table}

\begin{figure}[H]
\centering
\includegraphics[width=\textwidth]{fig_horizon.pdf}
\caption{Model performance versus prediction horizon. Left: ROC-AUC. Right:
PR-AUC, with the dotted line marking the no-skill base rate (which itself
falls as the horizon lengthens because the bankruptcy rate is lower in those
files). The tree ensembles preserve their advantage across all horizons; the
classical econometric models collapse toward random ranking on PR-AUC by H4.}
\label{fig:horizon}
\end{figure}

Three patterns stand out.

\paragraph{The classical models fall apart at longer horizons; the trees do
not.} On AUC, the three econometric models drop from around 0.78 at H1 to
around 0.60-0.63 at H4, close to random guessing. XGBoost barely moves,
slipping from 0.96 at H1 to 0.93 at H4. Random Forest behaves similarly. In
relative terms, the AUC gap between XGBoost and the best econometric model at
each horizon widens from about 17 percentage points at H1 to about 29 at H4.

\paragraph{PR-AUC reveals the collapse even more sharply.} Altman, Ohlson,
and Zmijewski drop to PR-AUC of 0.06-0.07 at H4, essentially the no-skill
base rate of 0.039: they no longer sort firms meaningfully at all. XGBoost
holds PR-AUC near 0.69, an eighteenfold lift over chance. This shows clearly
why simple accounting-ratio models are not reliable tools for medium-term
distress forecasting.

\paragraph{H5 partly recovers; H4 is the hardest horizon.} The non-monotonic
pattern, worst at H4 with a partial recovery at H5, appears across all six
models, so it is a feature of the data rather than any particular method.
The H5 file has the smallest sample (7,027 firms) and the lowest bankruptcy
rate (3.86\%), but appears to contain a more selected group of firms whose
financials already show clear signs of distress long before failure.
This is a feature of how \cite{zieba2016} constructed the five files; the
overall takeaway, that performance degrades sharply over the medium horizon
and that the ML advantage holds throughout, is unchanged.

\clearpage
% =============================================================================
\section{A Pooled Cross-Horizon Model: Horizon Dummies, Slopes, and Macro}\label{sec:pooled}
% =============================================================================

The horizon analysis in the previous section ran a separate model on each of
the five files. This section instead pools all five files into one sample and uses the
horizon structure directly, through horizon dummies and through interactions
between the ratios and the horizon. This is the part of a panel design the
data allow, as explained in Section~\ref{sec:panel-limits}. It lets us ask
two questions the separate-horizon models cannot answer. First, how much of
the cross-horizon difference is just a shift in the baseline default rate,
as opposed to a change in how the ratios work? Second, does a ratio's
predictive content stay the same as the horizon lengthens, or does it fade?

The pooled sample stacks the five horizon files into 43,405 firm-horizon
observations, of which 2,091 (4.82\%) are bankrupt. All models in this section
are logistic regressions estimated on standardized accounting ratios, so the
coefficients are comparable in size. We use the Ohlson-style ratio set, which
is available in every file.

\subsection{Horizon intercept dummies}

The baseline default rate falls steadily with the horizon, from 6.94\% at H1
to 3.86\% at H5 (Table~\ref{tab:horizon-grid}). A pooled model with a dummy for
each horizon (H1 as the omitted base) captures this directly. All four dummies
are strongly significant ($p < 10^{-8}$). Relative to H1, the odds of
bankruptcy are lower at every longer horizon, with odds ratios of 0.67 at H2,
0.60 at H3, 0.51 at H4, and 0.61 at H5. This is the dummy-variable counterpart
of the falling base rate: a firm that looks identical on its ratios is less
likely to fail within one year when observed further from the eventual outcome,
simply because the outcome is further away. The dummies play the role of the
period effects in a hazard model, absorbing the horizon-level shift in baseline
risk.

\subsection{Horizon-level macroeconomic variables}

Because each horizon file corresponds to a different average position in the
macro cycle, we can attach the horizon-level macro averages from
Table~\ref{tab:polish-macro} to the pooled data and ask whether they add
explanatory power. They do not, once the horizon dummies are present. The
reason is mechanical but worth stating: each horizon has a single set of macro
values, so the macro variables are collinear with the horizon dummies. Adding macro variables to a model that already has horizon dummies produces
no improvement (AUC 0.71 in both cases). At the horizon level, macro
variables and horizon dummies describe the same variation; the dummies are
more flexible. A genuine, separable macro effect would require firm-level
calendar dates, which the data do not provide. Reporting this null result
explicitly is more honest than simply leaving macro out.

\subsection{Do the ratio effects change with horizon?}

The more interesting question is whether the financial ratios mean the same
thing at every horizon. We test this by interacting each ratio with the horizon
dummies, the slope-dummy idea from Section~\ref{sec:panel-limits}. A joint
likelihood-ratio test on the full set of interaction terms strongly rejects the
hypothesis that the ratio effects are constant across horizons
($\chi^2 = 115.3$, 12 degrees of freedom, $p < 10^{-18}$). The ratios do behave
differently at different distances from the event.

Figure~\ref{fig:slope-horizon} shows the pattern, using single-ratio models so
that each coefficient keeps its natural economic sign. Table~\ref{tab:slopes}
reports the same numbers. The picture is consistent across the four ratios:
each one is at its strongest, in absolute terms, at the one-year horizon, and
weakens as the horizon lengthens. Profitability is the clearest example. A
one-standard-deviation increase in net income to assets lowers the bankruptcy
log-odds by 0.96 at H1, but by only 0.31 at H4, a drop of more than two thirds.
Working capital behaves the same way, from $-0.74$ at H1 to $-0.41$ at H4.
Leverage stays positive at every horizon, as expected, but its coefficient also
shrinks from 0.67 at H1 to 0.37 at H4. The partial rebound at H5 mirrors the
data artifact noted in Section~\ref{sec:horizons} and affects all the ratios
together.

\begin{table}[H]
\centering
\caption{Standardized single-ratio logit coefficients on bankruptcy by
prediction horizon, from models that interact each ratio with the horizon
dummies. Each coefficient is the effect of a one-standard-deviation increase in
the ratio on the bankruptcy log-odds at that horizon.}
\label{tab:slopes}
\small
\begin{tabular}{lccccc}
\toprule
Ratio & H1 & H2 & H3 & H4 & H5 \\
\midrule
Leverage (TL/TA)          & $+0.670$ & $+0.440$ & $+0.405$ & $+0.365$ & $+0.598$ \\
Profitability (NI/TA)     & $-0.956$ & $-0.575$ & $-0.444$ & $-0.307$ & $-0.687$ \\
Working capital (WC/TA)   & $-0.742$ & $-0.463$ & $-0.413$ & $-0.411$ & $-0.607$ \\
Liquidity (CA/STL)        & $-0.370$ & $-0.257$ & $-0.362$ & $-0.171$ & $-0.305$ \\
\bottomrule
\end{tabular}
\end{table}

\begin{figure}[H]
\centering
\includegraphics[width=0.85\textwidth]{fig_slope_horizon.pdf}
\caption{How the effect of each financial ratio on bankruptcy risk changes with
the prediction horizon. Each line is the standardized logit coefficient from a
single-ratio model interacted with the horizon dummies. All four ratios are
strongest near the event (H1) and weaken as the horizon lengthens, with a
partial rebound at H5 that reflects the same data feature seen elsewhere in the
paper.}
\label{fig:slope-horizon}
\end{figure}

This finding has a clear economic reading and connects back to the horizon
degradation in Section~\ref{sec:horizons}. The classical models lose so much
accuracy at longer horizons not because the ratios become noisier in some
mechanical sense, but because their genuine link to bankruptcy weakens. A
firm's current profitability says a great deal about whether it will fail
within a year and much less about whether it will fail within four. This also
helps explain why the tree ensembles hold up better across horizons. By
combining many ratios and exploiting interactions, they can lean on whatever
signal remains at longer horizons, while a model built on a few ratios with
fixed coefficients has less to fall back on once the strongest near-term
signals have faded.

\clearpage
% =============================================================================
\section{Opening the Black Box: What Drives the XGBoost Predictions}\label{sec:shap}
% =============================================================================

XGBoost is the most accurate model, but unlike the logit it does not
produce a simple table of coefficients. A natural concern is that its
advantage comes from fitting noise or variables with no economic meaning.
This section addresses that concern using SHAP values \citep{lundberg2017}.

\subsection{What SHAP values measure}

For each firm and each feature, a SHAP value measures how much that feature
pushed the firm's predicted log-odds of bankruptcy up or down, relative to
the average. The values come from cooperative game theory: the prediction is
treated as a payout that is fairly split among the contributing features.
For any firm, the SHAP values sum exactly to the gap between that firm's
prediction and the overall average, so they give a complete account of each
decision. Averaging absolute SHAP values across firms gives a global
importance ranking; the spread of SHAP values for a single feature shows the
direction and shape of its effect.

\subsection{Which features matter most}

Figure~\ref{fig:shap-bar} ranks the fifteen most important features by their
mean absolute SHAP value. Two things stand out. First, the variables XGBoost
relies on are economically sensible: operating-expense-to-liabilities
(\texttt{OpEx\_TL}), profit margins (\texttt{profit\_sales},
\texttt{profit\_sales\_TA}), sales growth, and several profitability ratios
(\texttt{EBITDA\_TA}, \texttt{NP\_DA\_TL}). These are exactly the kinds of
operating and profitability measures that the classical models also rely on,
which is reassuring, the model is not fitting noise. Second, two of the
top seven features are missingness flags (\texttt{miss\_OP\_finexp} and
\texttt{miss\_sales\_growth}), confirming once again that the pattern of
non-disclosure carries real information.

\begin{figure}[H]
\centering
\includegraphics[width=0.78\textwidth]{fig_shap_bar.pdf}
\caption{XGBoost global feature importance at horizon H1, measured by mean
absolute SHAP value. Blue bars are financial ratios; red bars are missingness
flags. The model relies mainly on profitability and cost ratios, with two
missingness flags among the most important features.}
\label{fig:shap-bar}
\end{figure}

Comparing this to the LASSO: in the linear LASSO, the missingness flags had
the \emph{largest} coefficients of all; in XGBoost, they remain important
but the financial ratios move to the top. The reason is that the tree
ensemble can extract more signal from the ratios directly, through nonlinear
thresholds and interactions, so it relies less on the missingness flags.

\subsection{The direction and shape of each effect}

Figure~\ref{fig:shap-beeswarm} is a beeswarm plot, which shows more than the
ranking. Each dot is one firm; its horizontal position is that firm's SHAP
value for the feature in that row, and its color is the value of the feature
(red for high, blue for low). This reveals the direction of each effect.

The two missingness flags show the clearest pattern: a tight cluster of red
dots (firms where the value \emph{is} missing, so the flag equals one)
sitting far to the right, meaning missingness sharply raises the predicted
risk. For the financial ratios the pattern is the expected one, for cost
and leverage measures, high values (red) tend to sit on the right
(higher risk), while for profitability measures high values tend to sit on
the left (lower risk). The fact that these directions line up with basic
financial intuition is further evidence that XGBoost's accuracy rests on a
sensible reading of the data rather than on artefacts.

\begin{figure}[H]
\centering
\includegraphics[width=0.85\textwidth]{fig_shap_beeswarm.pdf}
\caption{Per-firm SHAP values for the fifteen most important features at
horizon H1. Each dot is a firm; horizontal position is the feature's effect
on that firm's predicted log-odds; color is the feature value (red high, blue
low). The missingness flags (\texttt{miss\_OP\_finexp},
\texttt{miss\_sales\_growth}) show a clear cluster of high-value firms at far
right, meaning a missing value sharply raises predicted risk.}
\label{fig:shap-beeswarm}
\end{figure}

The SHAP analysis resolves the interpretability concern. XGBoost's most
influential variables are the operating and profitability measures that
decades of bankruptcy research has emphasized, the directions match
financial intuition, and the predictive power of missing values shows up
consistently across both the linear and nonlinear models.

\clearpage
% =============================================================================
\section{Robustness Checks}\label{sec:robustness}
% =============================================================================

The main results depend on three modeling choices: median imputation of
missing ratios, class weighting rather than resampling, and the use of the
full 75-feature set. This section stress-tests each choice with a targeted
check at horizon H1.

\subsection{Complete-case analysis}

Our headline numbers impute missing ratios and add missingness flags. To check
that this is not creating the results artificially, we re-run the analysis on
only the rows with no missing values at all (3{,}031 of 5{,}910 firms). This
complete-case sample is not a random subset: it excludes the distressed firms
whose missing values were themselves a signal, so its bankruptcy rate falls to
3.4\% and we should expect the machine-learning edge to shrink.

Table~\ref{tab:robust-cc} confirms this. XGBoost's AUC falls from 0.960 on the
full sample to 0.897 on complete cases, and its PR-AUC from 0.810 to 0.504.
The Ohlson logit, by contrast, holds up well and even improves slightly
(AUC 0.866). The interpretation is consistent with the rest of the paper: a
large part of XGBoost's advantage comes from exploiting the informative
missingness, so when those firms are removed the gap narrows. Even so, on the
complete-case sample XGBoost still clearly beats Ohlson on the rare-event
metric (PR-AUC 0.504 versus 0.244), so the conclusion that machine learning
helps is not an artefact of the imputation.

\begin{table}[H]
\centering
\caption{Complete-case analysis at H1 (rows with no missing values,
$n=3{,}031$, bankruptcy rate 3.4\%). Full-sample figures shown for comparison.}
\label{tab:robust-cc}
\small
\begin{tabular}{lcccc}
\toprule
Model & AUC & PR-AUC & Brier & Precision \\
\midrule
Ohlson (complete-case)  & 0.866 & 0.244 & 0.153 & 0.332 \\
XGBoost (complete-case) & 0.897 & 0.504 & 0.026 & 0.598 \\
\midrule
Ohlson (full sample)    & 0.795 & 0.297 & 0.176 & 0.324 \\
XGBoost (full sample)   & 0.960 & 0.810 & 0.030 & 0.830 \\
\bottomrule
\end{tabular}
\caption*{\footnotesize Source: author's calculations (code in Annex 1).}
\end{table}

\subsection{SMOTE instead of class weighting}

We handle the class imbalance by weighting the minority class in the loss
function. A common alternative is SMOTE, which creates synthetic bankrupt
firms to balance the training data \citep{chawla2002}. We re-run XGBoost with
SMOTE applied inside each training fold (never touching the validation fold)
and compare it to our class-weighted version. SMOTE gives slightly
\emph{lower} discrimination (AUC 0.937 versus 0.960, PR-AUC 0.724 versus
0.810) and a worse Brier score (0.044 versus 0.030). Class weighting is
therefore the better choice on this dataset, and it has the additional
advantage of not fabricating synthetic firms with financial-ratio
combinations that may never occur in reality. This also supports the broader
methodological point from Section~\ref{sec:lit}: resampling is not a free
lunch, and it can quietly damage calibration.

\subsection{Disentangling feature breadth from algorithmic power}
\label{sec:feature-breadth}

The previous tables show that XGBoost strongly outperforms the classical
benchmarks. That statement is true, but it is incomplete. XGBoost is not only
a different algorithm; it is also allowed to use a much wider set of predictors
than the textbook Altman, Ohlson, and Zmijewski models. A fair interpretation
therefore has to separate two channels: \emph{feature breadth}, meaning access
to more ratios and missingness indicators, and \emph{algorithmic flexibility},
meaning nonlinear thresholds and interactions.

There is one technical point worth being strict about. The labels
``Altman'', ``Ohlson'', and ``Zmijewski'' are not portable algorithms once the
feature set is changed; they are specific econometric specifications defined
by particular variables. Calling a 75-variable model ``Altman'' would be
misleading. The companion notebook therefore extends the robustness analysis
using a common class-weighted logit as the linear econometric algorithm and
then compares it with LASSO, random forest, and XGBoost on four feature sets:
Altman-5, Ohlson-7, the top-15 SHAP features, and all 75 features. This is the
clean decomposition: logit on 75 versus logit on 5 isolates feature breadth;
XGBoost on 75 versus logit on 75 isolates nonlinear modeling.

The XGBoost slice of this decomposition is already highly informative and is
reported in Table~\ref{tab:robust-features}. Restricting XGBoost to the five
Altman ratios or the seven Ohlson variables caps its AUC at about 0.82. Moving
to the top-15 SHAP variables raises AUC to 0.953 and PR-AUC to 0.788, already
close to the full 75-feature model. This means that the gain is not produced
by applying a sophisticated algorithm to the same handful of textbook ratios.
It comes mainly from letting the model see a broader set of financial signals,
including the missingness flags.

\begin{table}[H]
\centering
\caption{XGBoost performance by feature-set size at H1. The expanded notebook
runs the same feature-set grid for the transferable model classes (logit,
LASSO, random forest, XGBoost) and exports the full AUC and PR-AUC matrices.
The top-15 set is the SHAP importance ranking from Section~\ref{sec:shap}.}
\label{tab:robust-features}
\small
\begin{tabular}{lccc}
\toprule
Feature set & Number of features & AUC & PR-AUC \\
\midrule
Altman 5         &  5 & 0.816 & 0.336 \\
Ohlson 7         &  7 & 0.821 & 0.343 \\
Top 15 (by SHAP) & 15 & 0.953 & 0.788 \\
All 75           & 75 & 0.960 & 0.810 \\
\bottomrule
\end{tabular}
\caption*{\footnotesize Source: author's calculations (code in Annex 1).}
\end{table}

The full picture is clearest when the same feature sets are run across all the
model classes. Table~\ref{tab:feature-grid} reports this
grid and Figure~\ref{fig:featuregrid} plots it (for readability the figure
shows the four main classes; the regularized linear models are discussed
below). Reading down a column shows the
algorithm effect at a fixed information set; reading across a row shows the
feature-breadth effect for a fixed algorithm. Two regularities stand out. On a
narrow feature set (Altman-5 or Ohlson-7) the four algorithms are close
together, with AUC clustered around 0.77--0.82, because there is little for a
flexible model to exploit. Once the feature set widens to the top-15 or all-75
sets, every model improves and the tree ensembles pull clearly ahead of the
linear models. Feature breadth lifts all models; nonlinear modeling adds a
further premium only after the broader information set is available.

\begin{table}[H]
\centering
\caption{Performance by model class and feature set at H1.\textsuperscript{*}}
\label{tab:feature-grid}
\small
\begin{tabular}{lcccc}
\toprule
 & Altman-5 & Ohlson-7 & Top-15 & All-75 \\
\midrule
\multicolumn{5}{l}{\textit{Panel A: ROC-AUC}} \\
\midrule
Logit          & 0.772 & 0.795 & 0.905 & 0.892 \\
LASSO logit    & ---   & ---   & 0.907 & 0.895 \\
Ridge logit    & ---   & ---   & 0.905 & 0.877 \\
PCA logit      & ---   & ---   & ---   & 0.890 \\
Random Forest  & 0.818 & 0.804 & 0.950 & 0.938 \\
XGBoost        & 0.819 & 0.821 & 0.954 & 0.960 \\
\midrule
\multicolumn{5}{l}{\textit{Panel B: PR-AUC}} \\
\midrule
Logit          & 0.265 & 0.297 & 0.522 & 0.481 \\
LASSO logit    & ---   & ---   & 0.519 & 0.514 \\
Ridge logit    & ---   & ---   & 0.518 & 0.448 \\
PCA logit      & ---   & ---   & ---   & 0.521 \\
Random Forest  & 0.322 & 0.304 & 0.757 & 0.627 \\
XGBoost        & 0.335 & 0.338 & 0.789 & 0.810 \\
\bottomrule
\end{tabular}
\caption*{\footnotesize \textsuperscript{*} LASSO and Ridge are omitted on the Altman-5 and
Ohlson-7 sets because a penalty has little role with so few predictors. PCA
logit is reported only on the full 75-feature set, since dimensionality
reduction is meaningful only there; it uses the principal components explaining
95\% of the variance (about 30 components), fit inside each training fold. The
logit row is a plain class-weighted logistic regression, not the named
Altman/Ohlson specifications. Source: author's calculations (code in Annex 1).}
\end{table}

\begin{figure}[H]
\centering
\includegraphics[width=\textwidth]{fig_featuregrid.pdf}
\caption{Decomposing the machine-learning advantage at H1. Left: ROC-AUC.
Right: PR-AUC. Each line is one model class evaluated on feature sets of
increasing breadth. All models improve sharply when the feature set widens,
and the gap between the tree ensembles and the linear models opens up only at
the broader feature sets, the visual signature of feature breadth and
nonlinearity acting as two separate channels.}
\label{fig:featuregrid}
\end{figure}

A natural objection to the comparison so far is that a linear model with 75
correlated predictors is poorly specified: the financial ratios overlap
heavily, and a logit estimated on all of them at once is unstable, which could
understate what a linear model can really do. A standard econometric remedy is
principal component analysis, which replaces the correlated predictors with a
smaller set of uncorrelated components before estimation. To address this
directly, we add a PCA logit: within each training fold we standardize the 75
features, extract the principal components that explain 95\% of their variance
(about 30 components), and fit the logit on those components. The PCA is fit on
the training fold only, so the procedure stays leakage-free.

This is arguably the strongest case a linear model can make on this dataset:
it uses all of the information, and multicollinearity has been removed by
construction. It does improve on the raw 75-feature logit, reaching an AUC of
0.890 and a PR-AUC of 0.521 at H1, the best PR-AUC of any linear model here.
The result is robust to the variance threshold: keeping 90\%, 95\%, or 99\% of
the variance gives AUC of 0.880, 0.890, and 0.883 respectively. But it still
falls well short of XGBoost (AUC 0.960, PR-AUC 0.810). Removing
multicollinearity through PCA, in other words, does not close the gap. This
reinforces the central decomposition: once a linear model already has the full
feature set, neither variable selection (LASSO), coefficient shrinkage (Ridge),
nor decorrelation (PCA) lets it match the tree ensembles. What separates them
is nonlinearity, the ability to model thresholds and interactions among the
ratios, not the number or correlation structure of the inputs.

Table~\ref{tab:feature-decomp} gives the corresponding interpretation using
the already available H1 numbers. The narrow-feature algorithm gain is modest:
XGBoost on the Ohlson feature set improves AUC by only 0.026 relative to the
Ohlson logit. By contrast, moving from the Ohlson variables to the full
75-feature LASSO raises AUC by 0.100 and PR-AUC by 0.217 even though the model
remains linear. Once the broad feature set is available, XGBoost adds a further
0.065 AUC and 0.296 PR-AUC over the LASSO. The conclusion is not that
nonlinearity is irrelevant; it clearly matters. The conclusion is sharper:
feature breadth is a necessary condition for the ML advantage, and nonlinear
modeling extracts additional value only after the richer information set is
available.

\begin{table}[H]
\centering
\caption{A simple decomposition of the H1 performance gap. Positive numbers
are improvements in AUC or PR-AUC between the two models in the comparison.}
\label{tab:feature-decomp}
\small
\resizebox{\textwidth}{!}{%
\begin{tabular}{lccp{6.4cm}}
\toprule
Comparison & AUC gain & PR-AUC gain & Interpretation \\
\midrule
Ohlson logit (7) $\rightarrow$ XGBoost (Ohlson 7) & 0.026 & 0.046 & Nonlinear modeling on a narrow classical feature set adds only a small gain. \\
Ohlson logit (7) $\rightarrow$ LASSO logit (75) & 0.100 & 0.217 & A broad feature set substantially improves a still-linear model. \\
LASSO logit (75) $\rightarrow$ XGBoost (75) & 0.065 & 0.296 & Nonlinear thresholds and interactions add further value once the broad feature set is available. \\
XGBoost (Altman 5) $\rightarrow$ XGBoost (75) & 0.144 & 0.474 & Holding the algorithm fixed, feature breadth produces a very large gain. \\
\bottomrule
\end{tabular}%
}
\caption*{\footnotesize Source: author's calculations (code in Annex 1).}
\end{table}

This supports the central interpretation of the paper. The ML advantage is
real, but it should not be described as if the algorithm alone did the work.
It is the combination of a broader accounting-ratio set, explicit missingness
signals, and nonlinear tree ensembles that produces the large gap. This is
consistent with \cite{barboza2017}, who also find that adding complementary
financial indicators materially improves bankruptcy prediction across model
classes.


\clearpage
% =============================================================================
\section{Conclusion}\label{sec:conclusion}
% =============================================================================

This paper compared three classical bankruptcy models with modern
machine-learning classifiers on the Polish Companies Bankruptcy dataset,
using one evaluation protocol applied to every model. Keeping the protocol
fixed means that performance differences reflect the models themselves,
not differences in how they were trained or tested. Five findings emerge.

First, the machine-learning advantage on this dataset is large, not marginal.
At the one-year horizon, XGBoost reaches an AUC of 0.96 and a PR-AUC of 0.81,
against 0.80 and 0.30 for the best econometric model (Ohlson). The gap is much
wider on PR-AUC, the metric that matters most for a rare event, where XGBoost
delivers a nearly twelvefold lift over chance.

Second, the advantage is concentrated where it counts. At a common threshold,
XGBoost flags distressed firms with 83\% precision and 68\% recall, while
Ohlson manages 32\% and 46\%. The trees do not just improve the ranking;
they sharply reduce false alarms while still catching most failures.

Third, calibration and discrimination are separate properties, and conflating
them is a mistake. The class-weighted linear models are badly miscalibrated
out of the box but rank firms reasonably; the tree ensembles do both well. The
Brier-score decomposition shows that post-hoc recalibration removes almost all
of the calibration error without touching a model's underlying sorting power
and on that irreducible sorting power, the tree advantage is permanent.

Fourth, the advantage holds across horizons, and the pooled cross-horizon
model explains why. As the prediction horizon lengthens from one to five years,
the classical models decay toward random ranking on the rare-event metric,
while the tree ensembles retain most of their power. The pooled model shows
why this happens. When we let each ratio's effect vary by horizon, the effect
of profitability, working capital, and leverage is strongest at the one-year
horizon and weakens substantially further out. This pattern is jointly
significant at $p < 10^{-18}$. Simple accounting-ratio models are therefore
particularly ill-suited to medium-term distress forecasting, while flexible
models that combine many ratios cope better with the weaker long-horizon
signal.

Fifth, missing values are informative. Firms close to bankruptcy
systematically fail to report certain items, and the non-disclosure pattern
is among the strongest predictors in the data: the largest LASSO
coefficients and high on the XGBoost SHAP ranking. The robustness checks
confirm that a large part of the ML edge comes from this signal, which the
classical models cannot use.


Taken together, these findings also clarify the paper's contribution. The
missingness result is the main empirical addition: non-disclosure behaves like
a distress signal rather than a random data defect. The feature-breadth
analysis shows that the ML advantage should be decomposed rather than treated
as a black-box algorithmic victory. Finally, the unified protocol makes the
comparison credible, because the models differ in their structure and feature
sets but not in the rules under which they are trained and evaluated.

These findings come with limitations. Because firms are anonymized, we cannot
follow them over time, which rules out a proper hazard model and two Ohlson
variables that need prior-year data. The data cover Polish firms over
2000-2013, so the exact magnitudes may not apply elsewhere, though the
qualitative pattern, a large, robust ML advantage, likely holds more
broadly. We also used only structured financial ratios; adding market prices
or text disclosures could shift the balance between model types. The practical
case for any model also depends on how errors are valued; we used one
cost ratio, and a deployment decision would need to calibrate that
to the user's own context.

The broader message is methodological. The widely repeated claim that machine
learning beats classical credit-scoring models is, on this dataset and under a
fair protocol, correct. But the more useful results are the ones that a single
accuracy number hides: where the advantage comes from, whether the
probabilities can be trusted, how performance decays with the forecast horizon,
and which inputs are doing the work. Reporting those is what turns a model
comparison into evidence a practitioner can actually use.

% =============================================================================
% Bibliography
% =============================================================================

\bibliographystyle{abbrvnat}
\clearpage
\begin{thebibliography}{99}

\bibitem[Altman(1968)]{altman1968}
Altman, E.\ I.\ (1968).
\newblock Financial ratios, discriminant analysis and the prediction of
corporate bankruptcy.
\newblock \textit{Journal of Finance}, 23(4), 589-609.

\bibitem[Altman et~al.(1994)]{altman1994}
Altman, E.\ I., Marco, G., and Varetto, F.\ (1994).
\newblock Corporate distress diagnosis: comparisons using linear discriminant
analysis and neural networks (the Italian experience).
\newblock \textit{Journal of Banking \& Finance}, 18(3), 505-529.

\bibitem[Altman(2000)]{altman2000}
Altman, E.\ I.\ (2000).
\newblock Predicting financial distress of companies: revisiting the Z-score
and ZETA models.
\newblock Working paper, Stern School of Business, New York University.

\bibitem[Altman and Sabato(2007)]{altman2007}
Altman, E.\ I.\ and Sabato, G.\ (2007).
\newblock Modelling credit risk for SMEs: evidence from the US market.
\newblock \textit{Abacus}, 43(3), 332-357.

\bibitem[Barboza et~al.(2017)]{barboza2017}
Barboza, F., Kimura, H., and Altman, E.\ (2017).
\newblock Machine learning models and bankruptcy prediction.
\newblock \textit{Expert Systems with Applications}, 83, 405-417.

\bibitem[Breiman(2001)]{breiman2001}
Breiman, L.\ (2001).
\newblock Random forests.
\newblock \textit{Machine Learning}, 45(1), 5-32.

\bibitem[Brier(1950)]{brier1950}
Brier, G.\ W.\ (1950).
\newblock Verification of forecasts expressed in terms of probability.
\newblock \textit{Monthly Weather Review}, 78(1), 1-3.

\bibitem[Campbell et~al.(2008)]{campbell2008}
Campbell, J.\ Y., Hilscher, J., and Szilagyi, J.\ (2008).
\newblock In search of distress risk.
\newblock \textit{Journal of Finance}, 63(6), 2899-2939.

\bibitem[Carmona et~al.(2019)]{carmona2019}
Carmona, P., Climent, F., and Momparler, A.\ (2019).
\newblock Predicting failure in the U.S.\ banking sector: an extreme gradient
boosting approach.
\newblock \textit{International Review of Economics and Finance}, 61, 304-323.

\bibitem[Chawla et~al.(2002)]{chawla2002}
Chawla, N.\ V., Bowyer, K.\ W., Hall, L.\ O., and Kegelmeyer, W.\ P.\ (2002).
\newblock SMOTE: synthetic minority over-sampling technique.
\newblock \textit{Journal of Artificial Intelligence Research}, 16, 321-357.

\bibitem[Cohen(1988)]{cohen1988}
Cohen, J.\ (1988).
\newblock \textit{Statistical Power Analysis for the Behavioral Sciences},
2nd edition.
\newblock Lawrence Erlbaum Associates, Hillsdale, NJ.

\bibitem[Duan et~al.(2012)]{duan2012}
Duan, J.\ C., Sun, J., and Wang, T.\ (2012).
\newblock Multiperiod corporate default prediction, a forward intensity
approach.
\newblock \textit{Journal of Econometrics}, 170(1), 191-209.

\bibitem[Fisher(1936)]{fisher1936}
Fisher, R.\ A.\ (1936).
\newblock The use of multiple measurements in taxonomic problems.
\newblock \textit{Annals of Eugenics}, 7(2), 179-188.

\bibitem[Friedman(2001)]{friedman2001}
Friedman, J.\ H.\ (2001).
\newblock Greedy function approximation: a gradient boosting machine.
\newblock \textit{Annals of Statistics}, 29(5), 1189-1232.

\bibitem[Guo et~al.(2017)]{guo2017}
Guo, C., Pleiss, G., Sun, Y., and Weinberger, K.\ Q.\ (2017).
\newblock On calibration of modern neural networks.
\newblock \textit{Proceedings of the 34th International Conference on Machine
Learning}, PMLR 70, 1321-1330.

\bibitem[King and Zeng(2001)]{king2001}
King, G.\ and Zeng, L.\ (2001).
\newblock Logistic regression in rare events data.
\newblock \textit{Political Analysis}, 9(2), 137-163.

\bibitem[Lundberg and Lee(2017)]{lundberg2017}
Lundberg, S.\ M.\ and Lee, S.-I.\ (2017).
\newblock A unified approach to interpreting model predictions.
\newblock \textit{Advances in Neural Information Processing Systems}, 30,
4765-4774.

\bibitem[Mai et~al.(2019)]{mai2019}
Mai, F., Tian, S., Lee, C., and Ma, L.\ (2019).
\newblock Deep learning models for bankruptcy prediction using textual
disclosures.
\newblock \textit{European Journal of Operational Research}, 274(2), 743-758.

\bibitem[Min and Lee(2005)]{min2005}
Min, J.\ H.\ and Lee, Y.\ C.\ (2005).
\newblock Bankruptcy prediction using support vector machine with optimal
choice of kernel function parameters.
\newblock \textit{Expert Systems with Applications}, 28(4), 603-614.

\bibitem[Merton(1974)]{merton1974}
Merton, R.\ C.\ (1974).
\newblock On the pricing of corporate debt: the risk structure of interest rates.
\newblock \textit{Journal of Finance}, 29(2), 449-470.

\bibitem[Mossman et~al.(1998)]{mossman1998}
Mossman, C.\ E., Bell, G.\ G., Swartz, L.\ M., and Turtle, H.\ (1998).
\newblock An empirical comparison of bankruptcy models.
\newblock \textit{Financial Review}, 33(2), 35-54.

\bibitem[Murphy(1973)]{murphy1973}
Murphy, A.\ H.\ (1973).
\newblock A new vector partition of the probability score.
\newblock \textit{Journal of Applied Meteorology}, 12(4), 595-600.

\bibitem[Niculescu-Mizil and Caruana(2005)]{niculescu2005}
Niculescu-Mizil, A.\ and Caruana, R.\ (2005).
\newblock Predicting good probabilities with supervised learning.
\newblock \textit{Proceedings of the 22nd International Conference on Machine
Learning}, 625-632.

\bibitem[Odom and Sharda(1990)]{odom1990}
Odom, M.\ D.\ and Sharda, R.\ (1990).
\newblock A neural network model for bankruptcy prediction.
\newblock \textit{Proceedings of the International Joint Conference on Neural
Networks}, 2, 163-168.

\bibitem[Ohlson(1980)]{ohlson1980}
Ohlson, J.\ A.\ (1980).
\newblock Financial ratios and the probabilistic prediction of bankruptcy.
\newblock \textit{Journal of Accounting Research}, 18(1), 109-131.

\bibitem[Pellegrino et~al.(2022)]{pellegrino2022}
Lombardo, G., Pellegrino, M., Adosoglou, G., Cagnoni, S., Pardalos, P.\ M.,
and Poggi, A.\ (2022).
\newblock Machine learning for bankruptcy prediction in the American stock
market: dataset and benchmarks.
\newblock \textit{Future Internet}, 14(8), 244.

\bibitem[Pellegrino et~al.(2024)]{pellegrino2024}
Pellegrino, M., Lombardo, G., Adosoglou, G., Cagnoni, S., Pardalos, P.\ M.,
and Poggi, A.\ (2024).
\newblock A multi-head LSTM architecture for bankruptcy prediction with
time-series accounting data.
\newblock \textit{Future Internet}, 16(3), 79.

\bibitem[Saito and Rehmsmeier(2015)]{saito2015}
Saito, T.\ and Rehmsmeier, M.\ (2015).
\newblock The precision-recall plot is more informative than the ROC plot
when evaluating binary classifiers on imbalanced datasets.
\newblock \textit{PLoS ONE}, 10(3), e0118432.

\bibitem[Saerens et~al.(2002)]{saerens2002}
Saerens, M., Latinne, P., and Decaestecker, C.\ (2002).
\newblock Adjusting the outputs of a classifier to new a priori probabilities:
a simple procedure.
\newblock \textit{Neural Computation}, 14(1), 21-41.

\bibitem[Shin et~al.(2005)]{shin2005}
Shin, K.\ S., Lee, T.\ S., and Kim, H.\ J.\ (2005).
\newblock An application of support vector machines in bankruptcy prediction
model.
\newblock \textit{Expert Systems with Applications}, 28(1), 127-135.

\bibitem[Shumway(2001)]{shumway2001}
Shumway, T.\ (2001).
\newblock Forecasting bankruptcy more accurately: a simple hazard model.
\newblock \textit{Journal of Business}, 74(1), 101-124.

\bibitem[Smiti and Soui(2020)]{smiti2020}
Smiti, S.\ and Soui, M.\ (2020).
\newblock Bankruptcy prediction using deep learning approach based on
borderline SMOTE.
\newblock \textit{Information Systems Frontiers}, 22(5), 1067-1083.

\bibitem[Welch(1947)]{welch1947}
Welch, B.\ L.\ (1947).
\newblock The generalization of `Student's' problem when several different
population variances are involved.
\newblock \textit{Biometrika}, 34(1-2), 28-35.

\bibitem[Zięba et~al.(2016)]{zieba2016}
Zięba, M., Tomczak, S.\ K., and Tomczak, J.\ M.\ (2016).
\newblock Ensemble boosted trees with synthetic features generation in
application to bankruptcy prediction.
\newblock \textit{Expert Systems with Applications}, 58, 93-101.


\bibitem[Chava and Jarrow(2004)]{chava2004}
Chava, S.\ and Jarrow, R.\ (2004).
\newblock Bankruptcy prediction with industry effects.
\newblock \textit{Review of Finance}, 8(4), 537-569.

\bibitem[Chen and Guestrin(2016)]{chen2016}
Chen, T.\ and Guestrin, C.\ (2016).
\newblock XGBoost: a scalable tree boosting system.
\newblock \textit{Proceedings of the 22nd ACM SIGKDD International Conference
on Knowledge Discovery and Data Mining}, 785-794.

\bibitem[Dasilas and Rigani(2024)]{dasilas2024}
Dasilas, A.\ and Rigani, A.\ (2024).
\newblock Machine learning techniques in bankruptcy prediction: a systematic
literature review.
\newblock \textit{Expert Systems with Applications}, 255, 124761.

\bibitem[Hand(2009)]{hand2009}
Hand, D.\ J.\ (2009).
\newblock Measuring classifier performance: a coherent alternative to the area
under the ROC curve.
\newblock \textit{Machine Learning}, 77(1), 103-123.

\bibitem[Jones(1996)]{jones1996}
Jones, M.\ P.\ (1996).
\newblock Indicator and stratification methods for missing explanatory variables
in multiple linear regression.
\newblock \textit{Journal of the American Statistical Association}, 91(433),
222-230.

\bibitem[Jones et~al.(2017)]{jones2017}
Jones, S., Johnstone, D., and Wilson, R.\ (2017).
\newblock Predicting corporate bankruptcy: an evaluation of alternative
statistical frameworks.
\newblock \textit{Journal of Business Finance \& Accounting}, 44(1-2), 3-34.

\bibitem[Kostrzewa et~al.(2025)]{kostrzewa2025}
Kostrzewa, M.\ et~al.\ (2025).
\newblock Are foundation models useful for bankruptcy prediction? Evidence from
V4FinBench.
\newblock arXiv preprint.

\bibitem[Little and Rubin(2019)]{little2019}
Little, R.\ J.\ A.\ and Rubin, D.\ B.\ (2019).
\newblock \textit{Statistical Analysis with Missing Data}, 3rd edition.
\newblock Wiley.

\bibitem[Chakraborty and Ranjan(2024)]{granularimputation2024}
Chakraborty, D.\ and Ranjan, R.\ (2024).
\newblock Missing data imputation with granular semantics and AI-driven
pipeline for bankruptcy prediction.
\newblock arXiv:2404.00013.

\bibitem[Rubin(1976)]{rubin1976}
Rubin, D.\ B.\ (1976).
\newblock Inference and missing data.
\newblock \textit{Biometrika}, 63(3), 581-592.

\bibitem[Zmijewski(1984)]{zmijewski1984}
Zmijewski, M.\ E.\ (1984).
\newblock Methodological issues related to the estimation of financial
distress prediction models.
\newblock \textit{Journal of Accounting Research}, 22 (Supplement), 59-82.

\end{thebibliography}

\clearpage
% =============================================================================
\appendix
\section*{Annex 1. Key code for the reported calculations}
\addcontentsline{toc}{section}{Annex 1. Key code for the reported calculations}
\label{sec:appendix-code}
% =============================================================================

This annex reproduces the core of the code behind the tables marked ``Source:
author's calculations.'' It is not the full pipeline, which loads the dataset,
builds all 75 features, and produces every figure, but the parts that define
how models are scored and how the less standard models (Ridge, PCA logit, and
the pooled cross-horizon model) are estimated. All code is Python, using
\texttt{scikit-learn}, \texttt{xgboost}, and \texttt{statsmodels}.

\begin{center}
\fbox{\parbox{0.88\textwidth}{\centering
The complete, runnable code (a Jupyter notebook reproducing every table and
figure), this paper's \LaTeX{} source, and all figures are available in the
project repository:\\[6pt]
{\small\href{https://github.com/ilyas1python/Modeling-the-Probability-of-Corporate-Default-Using-Econometric-and-Machine-Learning-Methods}%
{\texttt{github.com/ilyas1python/Modeling-the-Probability-of-Corporate-Default-\\
Using-Econometric-and-Machine-Learning-Methods}}}
}}
\end{center}

\medskip

\subsection*{A.1 Evaluation metrics}

The two non-standard metrics are the Kolmogorov--Smirnov statistic and the
expected calibration error. The Brier score is decomposed into reliability,
resolution, and uncertainty by binning predictions.

\begin{lstlisting}[style=python]
def ks_statistic(y_true, y_prob):
    """Largest gap between the true-positive and false-positive rates across
    thresholds. 0 = no separation, 1 = perfect separation."""
    fpr, tpr, _ = roc_curve(y_true, y_prob)
    return float(np.max(tpr - fpr))

def expected_calibration_error(y_true, y_prob, n_bins=10):
    """Bin predictions into n_bins equal-width bins; average the gap between
    the empirical default rate and the mean predicted probability, weighted by
    bin size. 0 = perfectly calibrated."""
    edges = np.linspace(0.0, 1.0, n_bins + 1)
    idx = np.clip(np.digitize(y_prob, edges[1:-1], right=False), 0, n_bins - 1)
    N, ece = len(y_true), 0.0
    for b in range(n_bins):
        m = idx == b
        nb = int(m.sum())
        if nb:
            ece += (nb / N) * abs(y_true[m].mean() - y_prob[m].mean())
    return float(ece)
\end{lstlisting}

\subsection*{A.2 Unified cross-validation protocol}

Every model in the paper is scored through one harness. It runs stratified
five-fold cross-validation, imputes missing values with the training-fold
median, optionally scales features (for linear models), optionally applies
isotonic recalibration on an inner holdout, and returns mean metrics across
folds. Fitting all preprocessing inside the training fold is what keeps the
comparison leakage-free.

\begin{lstlisting}[style=python]
def run_cv(clf, X, y, num_cols, passthrough_cols=(), scale=False,
           recalibrate=False, n_splits=5, seed=42, c_fn=35.0, c_fp=1.0):
    """Stratified k-fold CV for any scikit-learn classifier. Returns mean and
    std of every metric across folds. All preprocessing is fit on the training
    fold only (leakage-free)."""
    skf = StratifiedKFold(n_splits=n_splits, shuffle=True, random_state=seed)
    rows = []
    for tr, va in skf.split(X, y):
        # Median imputation fit on the training fold only.
        imp = SimpleImputer(strategy='median').fit(X.iloc[tr][num_cols])
        Xtr = np.hstack([imp.transform(X.iloc[tr][num_cols]),
                         X.iloc[tr][list(passthrough_cols)].values])
        Xva = np.hstack([imp.transform(X.iloc[va][num_cols]),
                         X.iloc[va][list(passthrough_cols)].values])
        if scale:                      # recommended for linear models
            sc = RobustScaler().fit(Xtr)
            Xtr, Xva = sc.transform(Xtr), sc.transform(Xva)
        model = clone(clf).fit(Xtr, y.iloc[tr])
        p = model.predict_proba(Xva)[:, 1]
        if recalibrate:                # isotonic fit on an inner 25% holdout
            p = _isotonic_recalibrate(model, Xtr, y.iloc[tr], Xva)
        rows.append(_all_metrics(y.iloc[va].values, p, c_fn, c_fp))
    return pd.DataFrame(rows).agg(['mean', 'std']).T
\end{lstlisting}

\subsection*{A.3 Econometric baselines and the seven models}

The classical models are class-weighted logits on their original variable
sets; the regularized linear models and tree ensembles use all 75 features.
This is the specification list used for the main comparison and the horizon
grid.

\begin{lstlisting}[style=python]
def _logit():
    return LogisticRegression(class_weight='balanced', max_iter=5000,
                              solver='lbfgs')

model_specs = [
    ('Altman (5)',       _logit(),        ALTMAN_5,    (),        True),
    ('Ohlson (7+OENEG)', _logit(),        OHLSON_7,    ['OENEG'], True),
    ('Zmijewski (3)',    _logit(),        ZMIJEWSKI_3, (),        True),
    ('LASSO logit (75)', LogisticRegression(penalty='l1', solver='liblinear',
        C=0.1, class_weight='balanced', max_iter=2000), FEATURES, MISS_INDS, True),
    ('Ridge logit (75)', LogisticRegression(penalty='l2', solver='lbfgs',
        C=1.0, class_weight='balanced', max_iter=2000), FEATURES, MISS_INDS, True),
    ('Random Forest',    RandomForestClassifier(n_estimators=300,
        max_features='sqrt', class_weight='balanced', n_jobs=2,
        random_state=42), FEATURES, MISS_INDS, False),
    ('XGBoost',          XGBClassifier(n_estimators=400, learning_rate=0.05,
        max_depth=4, subsample=0.8, colsample_bytree=0.8, reg_lambda=1.0,
        min_child_weight=5, scale_pos_weight=spw), FEATURES, MISS_INDS, False),
]

for label, clf, num_cols, pass_cols, scale in model_specs:
    summary = run_cv(clf, X=Xh1, y=Xh1['y'], num_cols=num_cols,
                     passthrough_cols=pass_cols, scale=scale)
\end{lstlisting}

\subsection*{A.4 PCA logit}

The PCA logit answers the objection that a logit on 75 correlated ratios is
poorly specified. Principal components are extracted inside each training fold
(keeping 95\% of variance) and the logit is fit on the components, so the
procedure stays leakage-free.

\begin{lstlisting}[style=python]
def pca_logit_cv(Xh, var_target=0.95, seed=42):
    y = Xh['y'].astype(int).values
    X = Xh[ALL_FEATURES].values
    skf = StratifiedKFold(n_splits=5, shuffle=True, random_state=seed)
    aucs, prs, ncomps = [], [], []
    for tr, va in skf.split(X, y):
        imp = SimpleImputer(strategy='median').fit(X[tr])
        Xtr, Xva = imp.transform(X[tr]), imp.transform(X[va])
        sc = StandardScaler().fit(Xtr)
        Xtr, Xva = sc.transform(Xtr), sc.transform(Xva)
        # PCA fit on the training fold only.
        pca = PCA(n_components=var_target, svd_solver='full').fit(Xtr)
        Ztr, Zva = pca.transform(Xtr), pca.transform(Xva)
        ncomps.append(Ztr.shape[1])
        m = LogisticRegression(class_weight='balanced', max_iter=2000,
                               random_state=seed).fit(Ztr, y[tr])
        p = m.predict_proba(Zva)[:, 1]
        aucs.append(roc_auc_score(y[va], p))
        prs.append(average_precision_score(y[va], p))
    return np.mean(aucs), np.mean(prs), int(np.mean(ncomps))
\end{lstlisting}

\subsection*{A.5 Pooled cross-horizon model}

The pooled model stacks the five horizon files, adds a dummy for each horizon,
and interacts the financial ratios with the horizon dummies. The likelihood-ratio
test compares the model with interactions against the model with horizon dummies
only, testing whether the ratio effects change with the horizon.

\begin{lstlisting}[style=python]
import statsmodels.api as sm
from scipy import stats

# Stack the five horizon files with horizon dummies.
pool = pd.concat([horizons[h].assign(horizon=h)
                  for h in ['H1', 'H2', 'H3', 'H4', 'H5']], ignore_index=True)
y = pool['y'].astype(int).values
Hd = pd.get_dummies(pool['horizon'], prefix='H', drop_first=True).astype(float)

# Model B: ratios + horizon intercept dummies.
XB = pd.concat([standardize(pool[OHLSON_7]), Hd], axis=1)
mB = sm.Logit(y, sm.add_constant(XB)).fit(disp=0)

# Model D: add ratio x horizon interactions (slope dummies).
key = ['TL_TA', 'net_profit_TA', 'WC_TA']
inter = {f'{k}_x_{c}': X0[k].values * Hd[c].values for k in key for c in Hd.columns}
XD = pd.concat([XB, pd.DataFrame(inter)], axis=1)
mD = sm.Logit(y, sm.add_constant(XD)).fit(disp=0)

# Likelihood-ratio test for the joint significance of the interactions.
LR = 2 * (mD.llf - mB.llf)
df = XD.shape[1] - XB.shape[1]
p_value = stats.chi2.sf(LR, df)        # chi2 = 115.3, df = 12, p < 1e-18
\end{lstlisting}


\end{document}
